\documentclass[11pt,a4paper]{amsart}
\usepackage{amssymb}
\usepackage{graphicx}
%\usepackage[dvips]{graphicx}
%\usepackage{psfrag}

%\newtheorem{Thm}{Theorem}[section]
%\newtheorem{Prop}[Thm]{Proposition}
%\newtheorem{Lem}[Thm]{Lemma}
%\newtheorem{Cor}[Thm]{Corollary}
%\newtheorem{Claim}{Claim}
\usepackage[colorlinks,
linkcolor=blue,
anchorcolor=blue,
citecolor=blue]{hyperref}
\newcommand{\eref}[1]{(\ref{#1})}

\def\cv{\ensuremath{\text {Cor}}}
\def\ld{\ensuremath{\text{LD}}}

\theoremstyle{plain}
\newtheorem{Main}{Theorem}
\renewcommand{\theMain}{\Alph{Main}}
\newtheorem{Thm}{Theorem}[section]
\newtheorem{Prop}[Thm]{Proposition}
\newtheorem{Lem}[Thm]{Lemma}
\newtheorem*{Sublem}{Sublemma}
\newtheorem*{Claim}{Claim}
\newtheorem{Cor}[Thm]{Corollary}
\newtheorem*{Fact}{Fact}

\theoremstyle{remark}
\newtheorem{Rem}[Thm]{Remark}
\newtheorem{Def}[Thm]{Definition}
\newtheorem{Exa}[Thm]{Example}
\newtheorem*{Prb}{Problem}
\newtheorem*{Not}{Notation}

\def\real{{\mathbb R}}
\def\rational{{\mathbb Q}}
\def\natural{{\mathbb N}}
\def\integer{{\mathbb Z}}
\def\B{\mathcal B}
\def\E{\mathcal{E}}
\def\H{\mathcal{H}}
\def\O{\mathcal{O}}
\def\P{\mathcal{P}}
\def\R{\mathcal{R}_{\epsilon,\de}}
\def\S{\mathcal{S}}
\def\Ga{\Gamma_{\epsilon,\de,n}}
\def\F{\tilde{F}_2}



\def\all{\operatorname{all}}
\def\const{\operatorname{const}}
\def\closure{\operatorname{closure}}
\def\dist{\operatorname{dist}}
\def\graph{\operatorname{graph}}
\def\inc{\operatorname{inc}}
\def\Int{\operatorname{Int}}
\def\id{\operatorname{id}}
\def\Lip{\operatorname{Lip}}
\def\max{\operatorname{max}}
\def\min{\operatorname{min}}
\def\mod{\operatorname{mod}}
\def\o{\operatorname{o}}
\def\rem{\operatorname{rem}}
\def\sign{\operatorname{sign}}
%\def\sup{\operatorname{sup}}
\def\supp{\operatorname{supp}}
\def\W{\operatorname{W}}

\def\al{\alpha}
\def\be{\beta}
\def\de{\delta}
\def\vep{\varepsilon}
\def\vap{\varphi}
\def\th{\theta}
\def\om{\omega}
\def\la{\lambda}
\def\La{\Lambda}
\def\ga{\gamma}
\def\ka{\kappa}
\def\pa{\partial}

\addtolength{\topmargin}{-0.5cm}
\addtolength{\textheight}{0cm}

\setlength{\textwidth}{14.3cm}
\setlength{\textheight}{23cm}
\addtolength{\oddsidemargin}{-0.82cm}
\addtolength{\evensidemargin}{-0.82cm}
\begin{document}
	
	\title[]
	{the dimension of the non-dense orbit set for expanding maps}
	\author{Congcong Qu}
	\address{Congcong Qu, College of Big Data and software Engineering, Zhejiang Wanli University, Ningbo, 315107, Zhejiang, P.R.China}
	\email{congcongqu@foxmail.com}
	\iffalse
	\author{Yongluo Cao \textsuperscript{*}}
	\address{Yongluo Cao, Department of Mathematics, Soochow University,
		Suzhou 215006, Jiangsu, P.R. China}
	\address{and Center for Dynamical Systems and Differential Equations, Soochow University, Suzhou 215006, Jiangsu, P.R. China}
	\email{ylcao@suda.edu.cn}
	
	\date{\today}
	
	
	\thanks{
		* Y. Cao is the corresponding author.
	}
	\fi
	\subjclass[2000]{37C45,37D05,37D20}
	
	\keywords{Hausdorff dimension, repeller, non-dense orbit}
	
	
	
	\begin{abstract}
		In this paper, we give a new proof for  the Hausdorff dimension of the non-dense orbit set for expanding maps. This proof is based on the sharp lower bound of the Hausdorff dimension of  repellers given by Cao, Pesin and Zhao in~\cite{caopesinzhao2019}. Besides, we prove that the Carath$\acute{\text{e}}$odory singular dimension of non-dense orbit set equals that of the repeller.
	\end{abstract}
	
	\maketitle
	
	%\tableofcontents
	
	\section{Introduction}
	Suppose ~$M$ is a compact~Riemannian manifold and~$f:M\rightarrow M$ is a continuous map. Given an~$f-$invariant set~$\Lambda\subset M$ and a subset~$A\subset \Lambda$, the ~$A-$exceptional set in~$\Lambda$ with respect to $f|_{\Lambda}$ is defined as~$E^{+}_{f|_{\Lambda}}(A)=\{x\in \Lambda:\overline{\mathcal{O}_{f}(x)}\cap A={\varnothing}\}$, where $\mathcal{O}_{f}(x)$ denotes the forward orbit of $x$ under the action of $f$. Denote by $ND(f)$ the set of points with non-dense forward orbit under $f$. By definition, $E^+_{f|_M}(y)\subset ND(f)$ for any $y\in M$. If $f$ preserves an ergodic measure of full support, then $ND(f)$ has measure $0$. It's natural to consider the dimension of the non-dense orbit set.
	
	In this paper, we consider the  dimension of the non-dense orbit set for $C^{1+\gamma}$ expanding maps. Precisely, we prove the following results. Here we use $E^{+}_{f}(y)$ shortly for $E^{+}_{f|_{M}}(y)$. We use $dim_C Z$ to denote Carath$\acute{\text{e}}$odory singular dimension of the set $Z$, see subsection \ref{C dimension} for definition.
	\begin{Main}\label{expanding}
		Let $M$ be a $d-$dimensional Riemannian manifold and $f: M\rightarrow M$ be a $C^{1+\gamma}$ expanding map. Suppose that $\mu$ is an ergodic measure for $f$ which is absolutely continuous to the Lebesgue measure. Then for any $y\in M$, we have $$dim_{H} E^{+}_{f}(y)=dim_{H}\mu=d.$$
	\end{Main}
	\begin{Main}\label{expanding2}
		Let $M$ be a $d-$dimensional Riemannian manifold and $f: M\rightarrow M$ be a $C^{1+\gamma}$ map admitting a repeller $\Lambda$. Then for any $y\in \Lambda$, we have $$dim_{C} E^{+}_{f|_{\Lambda}}(y)=dim_{C} \Lambda.$$
	\end{Main}
	\iffalse
	\begin{Main}\label{hyperbolic1}
		Let $M$ be a $d-$dimensional Riemannian manifold and $f: M\rightarrow M$ be a transitive $C^{1+\alpha}$ hyperbolic diffeomorphism. Suppose that $\mu$ is the unique SRB measure for $f$. Then for any $x, y\in M$, we have $$dim_{H} E^{+}_{f}(y)\cap W^{u}(x)=dim_{H}\mu_{x}^{u}=dim W^{u}(x),$$
		where $\mu_{x}^{u}$ is the conditional measure on the unstable leaf.
	\end{Main}
	\begin{Main}\label{hyperbolic2}
		Let $M$ be a $d-$dimensional Riemannian manifold and $f: M\rightarrow M$ be a $C^{1+\alpha}$ diffeomorphism. Suppose $\Lambda\subset M$ is a locally maximal hyperbolic invariant set. Then for any $x,y\in\Lambda$, we have $$dim_{C} E^{+}_{f|_{\Lambda}}(y)\cap \W^{u}(x)=dim_{C} \Lambda\cap W^{u}(x).$$
	\end{Main}
\fi
	For the history of the research on the dimension of the exceptional set, one can refer to ~\cite{CG1, CG2}. Now we state the differences of our results with the existing ones. The first result for the dimensions of the exceptional sets for expanding maps and hyperbolic diffeomorphsims was obtained by  Urb$\acute{\text{a}}$nski in \cite{Urbanski}. He showed the results of our Main Theorem ~\ref{expanding}  and for Main Theorem ~\ref{expanding2} in conformal case. Notice that for conformal maps, the Hausdorff dimension and the Carath$\acute{\text{e}}$odory singular dimension coincide. He showed the results by a good lower bound of the Hausdorff dimension of a suitable invariant measure supported on the manifold, i.e. volume measure for expanding maps and the Gibbs measure for Anosov diffeomorphisms respectively. And Tseng claimed in \cite{Tseng1} there is a minor gap in Urb$\acute{\text{a}}$nski's proof for the lower bound of the Hausdorff dimension of the measure, and gave a modification of that by playing the Schmidt game, which is developed by Schmidt in \cite{Schmidt}, for expanding endomorphisms on the circle in~\cite{Tseng1} and for certain Anosov diffeomorphisms on the two dimensional torus in ~\cite{Tseng2}. Subsequent work is due to Wu in ~\cite{Wu1, Wu2, Wu3}. His results are based on the tool of Schmidt game and the modified Schmidt game which was introduced by Kleinbock and Weiss in \cite{KW}. Wu extended Tseng's results to the general expanding endomorphisms and partially hyperbolic diffeomorphisms admitting unstable directions.
	
	
	
	Other interesting results are due to Campos and Gelfert in ~\cite{CG1,CG2}. In ~\cite{CG1} and ~\cite{CG2}, for conformal $C^{1+\gamma}$ endomorphisms and diffeomorphisms admitting a locally maximal invariant set, Campos and Gelfert proved that for a set with either small entropy or dimension, the exceptional set of such a set is large in the sense of the Hausdorff dimension and entropy. Their work was heavily based on Dolgopyat's result on entropy in \cite{Do} on the symbolic spaces. They constructed repellers and hyperbolic sets inside the locally maximal invariant set respectively to approximate the entropy and the Lyapunov exponents. The semi-conjugacy and conjugacy respectively of the repellers and hyperbolic sets to the symbolic systems yields the result for the entropy. And the relation of the Hausdorff dimension, entropy and the Lyapunov exponents by Young~\cite{Young} gives the results of dimension. The first author and Juan Wang extended these results to average conformal case in~\cite{QW}.
	
	Our proof of Theorem~\ref{expanding} is based on the lower bound of the Hausdorff dimension of a repeller obtained by Cao, Pesin and Zhao in \cite{caopesinzhao2019}. We briefly state our proof here. For an ergodic measure $\mu$ for $f$ which is absolutely continuous with the Lebesgue measure and a given point $y\in M$, we can construct a sequence of repellers $\{\Lambda_n\}_{n\in\mathbb{N}}$  to approximate the measure $\mu$ and  $\{\Lambda_n\}_{n\in\mathbb{N}}$ can be chosen to avoid the point $y$ and hence $\Lambda_n \subset E^{+}_{f}(y)$. Since $\mu$ is absolutely continuous with respect to the Lesbegue measure, the Pesin's formula holds. We can use this to deduce that the Hausdorff dimension of $\{\Lambda_n\}_{n\in\mathbb{N}}$ are close to the dimension of the manifold via the lower bound of the Hausdorff dimension of the repellers obtained in~\cite{caopesinzhao2019}. By the monotonicity of the Hausdorff dimension, we can show that $dim_H E^{+}_{f}(y)$, $dim_H \mu$ and the dimension of the manifold coincide.
	
	The results in Theorem~\ref{expanding2} is brand new. Cao, Pesin and Zhao \cite{caopesinzhao2019} obtained that the Carath$\acute{\text{e}}$odory singular dimension of the repeller is exactly the unique root of the Bowen's equation involving the sub-additive topological pressure of the sub-additive singular-valued potential. Thus for a given repeller $\Lambda$ and $y\in \Lambda$, we can find a sequence of repellers $\Lambda_n\subset \Lambda$ and $\Lambda_n\cap \{y\}=\varnothing$, and then use  the Carath$\acute{\text{e}}$odory singular dimension of $\Lambda_n$ to approximate that of $\Lambda$.
	\iffalse
	The idea of the proof of Main Theorem ~\ref{hyperbolic1} and ~\ref{hyperbolic2} are similar to the proof of the Main Theorem~\ref{expanding} and ~\ref{expanding2} respectively.
	\fi
	
	This paper is organized as follows. In section \ref{Preliminary}, we recall some basic definitions and preliminary results. In section \ref{proof}, we give the details of the proof of the main result.
	\section{Some definitions and preliminary results}\label{Preliminary}
	In this section, we recall some notions used in our proof and state some preliminary results.

	\subsection{Repellers for expanding maps.} Let $M$ be a $d-$dimensional compact $C^{\infty}$ Riemannian manifold  and $f:M\rightarrow M$ be a $C^{1}$ map. Suppose $\Lambda$ is a compact $f-$invariant subset of $M$. We say $\Lambda$ is {\it a repeller for $f$ and $f$ is expanding} if there exists an open neighborhood $U$ of $\Lambda$ such that $\Lambda=\{x\in U: f^{n}(x)\in U, \text{for~all}~n\in\mathbb{N}\}$ and there is $\lambda>0$ such that for all $x\in \Lambda$ and $v\in T_{x}M$, we have $\|D_{x}f (v)\|\geq \lambda \|v\|$, where $\|\cdot\|$ is the norm induced by the Riemannian metric on $M$.
	\iffalse
	Suppose $f$ is a $C^{1+\alpha}$ diffeomorphism on $M$ and $\Lambda\subset M$ is an $f-$invariant compact set. We say $\Lambda$ is a hyperbolic set if for any $x\in \Lambda$, the tangent space admits a decomposition $T_{x}M=E^{s}(x)\oplus E^{u}(x)$ with the following properties:
	\begin{enumerate}
		\item The splitting is $Df-$invariant: $D_{x}f E^{\sigma}(x)=E^{\sigma}(f(x))$ for $\sigma=s,u$.
		\item The stable subspace $E^{s}(x)$ is uniformly contracting and the unstable subspace $E^{u}(x)$ is uniformly expanding: there are constants $C\geq 1$ and $0<\chi<1$ such that for every $n\geq 0$ and $v^{s,u}\in E^{s,u}(x)$, we have $$\|D_{x}f^{n}v^{s}\|\leq C\chi^{n}\|v^{s}\|\quad \text{and} \quad \|D_{x}f^{-n}v^{u}\|\leq C\chi^{n}\|v^{u}\|.$$
	\end{enumerate}
	
	We say $\Lambda$ is locally maximal if there exists an open neighborhood $U$ of $\Lambda$ such that $\Lambda=\bigcap_{n\in\mathbb{Z}}f^{n}U$.
	\fi
	\iffalse
	\subsection{Markov partition and Gibbs measures}
	Assume that $\Lambda$ is a locally maximal hyperbolic set for a $C^{1}$ diffeomorphism $f$ and $f|_{\Lambda}$ is transitive. A finite cover $\mathcal{P}=\{P_{1},P_{2},...,P_{l}\}$ of $\Lambda$ is called a Markov partition of $\Lambda$ if
	\begin{enumerate}
		\item $P_{i}\neq \varnothing$ and $P_{i}=\overline{\text{int} P_{i}}$,
		\item $\text{int} P_{i} \cap \text{int} P_{j} =\varnothing$ whenever $i\neq j$,
		\item for each $x\in \text{int} P_{i}\cap f^{-1} \text{int} P_{j}$, we have
		$$f(W^{s}_{\text{loc}}(x)\cap P_{i})\subset W^{s}_{\text{loc}}(f(x))\cap P_{j},$$
		$$f(W^{u}_{\text{loc}}(x)\cap P_{i})\supset W^{u}_{\text{loc}}(f(x))\cap P_{j}.$$
	\end{enumerate}
	It is well-known that a locally maximal hyperbolic set $\Lambda$ admits Markov partitions of arbitrary small diameters. One can refer to \cite{KH}.
	
	Given a Markov partition $\{P_{1},P_{2},...,P_{l}\}$, a sequence $\mathbf{i}=(i_{-m}...i_{0}...i_{n})$, where $1\leq i_{j}\leq l$, is called admissible if for $j=-m,...,0,...,n-1$, $\text{int} P_{j}\cap f^{-1}\text{int} P_{j+1}\neq \varnothing$. Given an admissible sequence $\mathbf{i}=(i_{-m}...i_{0}...i_{n})$, one can define the cyclinder
	$$P_{i_{-m}...i_{0}...i_{n}}=\bigcap_{j=-m}^{n}f^{-j}P_{i_{j}}.$$
	
	A Markov partition $\{P_{1},P_{2},...,P_{l}\}$ generates a symbolic model of $\Lambda$ by a subshift of finite type $(\Sigma_{A},\sigma)$, where $\Sigma_{A}$ is the set of two-sided infinite sequences of integers which are admissible with respect to the transfer matrix of the Markov partition $A=(a_{ij})$ (i.e. $a_{ij}=1$ if $\text{int} P_{i}\cap f^{-1}\text{int} P_{j}\neq \varnothing$ and $a_{ij}=0$ otherwise).
	
	Now we can define the coding map $\pi:\Sigma_{A}\rightarrow \Lambda$ by
	\begin{align*}
		\pi(...i_{-n}...i_{0}...i_{n}...)=\bigcap_{n\geq 0}P_{i_{-n}...i_{n}}.
	\end{align*}
	
	We say a Borel probability measure $\mu$ on $\Lambda$ is a Gibbs measure for a continuous function $\varphi$ on $\Lambda$ if there exists $C>0$ such that for every $n\in \mathbb{N}$, every admissible sequence $(i_{0}i_{1}...i_{n-1})$, and $x\in P_{i_{0}i_{1}...i_{n-1}}$, we have
	\begin{align*}
		C^{-1}\leq \frac{\mu(P_{i_{0}i_{1}...i_{n-1}})}{\exp(-nP+S_{n}\varphi(x))}\leq C,
	\end{align*}
	where $P$ is a constant and $S_{n}\varphi(x)=\sum_{j=0}^{n-1}\varphi(f^{j}(x))$.
	\fi
	\iffalse
	\begin{Prop}
		Suppose $\Lambda$ is a locally maximal hyperbolic set for a $C^{1+\alpha}$ diffeomorphism and $f|_{\Lambda}$ is transitive. Assume $\varphi:\Lambda\rightarrow \mathbb{R}$ is a H$\ddot{o}$lder potential. Then there exists a unique equilibrium state $\mu$. And the family of the conditional measures $\{\mu_{x}^{u}\}_{x\in\Lambda}$ of $\mu$ on the unstable leaf has the $u-$Gibbs property for the potential $\varphi$, that is, there is $Q>0$ such that
		for every  $x\in\Lambda$, if $\pi(...i_{-m}...i_{0}...i_{n}...)=x$, then for each $n\in \mathbb{N}$ and $y\in P_{...i_{-m}...i_{0}...i_{n-1}}$, we have
		\begin{align*}
			Q^{-1}\leq \frac{\mu_{x}^{u}(P_{...i_{-m}...i_{0}...i_{n-1}})}{\exp(-nP(\varphi)+S_{n}\varphi(y))}\leq Q,
		\end{align*}
		where $P(\varphi)$ is the topological pressure.
	\end{Prop}
	\fi
	\subsection{Metric entropy, topological entropy and the classical topological pressure}
	For the definitions of the metric entropy, topological entropy and the classical topological pressure, one can refer to ~\cite{Pesin} or~\cite{WALTERS}. Here we recall the variational principle of the topological entropy.
	
	Consider a continuous transformation $f: X\to X$ of a compact metric space $X$ equipped with a metric $d$. Suppose that $\Lambda\subset X$ is an $f-$invariant compact set. For an invariant measure $\mu$ for $f$ which is supported on $\Lambda$, we denote by $h_{\mu}(f|_{\Lambda})$ the metric entropy and by $h_{\text{top}}(f|_{\Lambda})$ the topological entropy on $\Lambda$. Classical result shows the following variational principle holds. For the proof, one can refer to ~\cite{WALTERS}.
	\begin{Thm}[Variational principle]
		\begin{align*}
			h_{\text{top}}(f|_{\Lambda})=\sup\{h_{\mu}(f|_{\Lambda})|\supp(\mu)\subset\Lambda,~\mu~\text{is~ergodic}\}.
		\end{align*}
	\end{Thm}
	The classical result shows that for an expansive map, the map $\mu\mapsto h_{\mu}(f|_{\Lambda})$ is upper semi-continuous. For this, one can refer to ~\cite{WALTERS}. Notice that expanding maps  are expansive. Hence one can find an ergodic measure to achieve the supreme.
	\iffalse
	\subsection{Lyapunov exponents and the ~Oseledec's Theorem}
	Next we recall the definition of the Lyapunov exponents and the ~Oseledec's Theorem.
	
	\begin{Thm} \cite{Oseledets} Suppose $f$ is  a  $C^{1}$ map on a $d$-dimensional compact Riemannian manifold $M$ and $\mu$ is an  $f$-invariant measure, then there exists an $f$ invariant set  $R$, $(i.e.~ f(R)\subset R)$ with $\mu(R)=1$ such that for each $x\in R$,
		\begin{enumerate}
			\item  There exists a measurable filtration
			\begin{displaymath}
				\{0\}=E_{0}(x)\subset E_{1}(x)\subset...\subset E_{\ell(x)}(x)=T_{x}M,
			\end{displaymath}
			where $\{E_{i}(x)\}_{1\leq i\leq \ell(x)}$ are linear subspaces of $T_{x}M$ such that
			\begin{displaymath}
				E_{i}(x)\neq E_{i+1}(x)\quad \text{and} \quad D_{x}f E_{i}(x)=E_{i}(f(x)).
			\end{displaymath}
			\item And there exist measurable functions $\lambda_{1}(x)<\lambda_{2}(x)<...<\lambda_{\ell(x)}(x)$, such that for any $v\in E_{i}(x)\setminus E_{i-1}(x), 1\leq i\leq \ell(x)$, we have~$$\lambda_{i}(x)=\lim_{n\rightarrow \infty}\frac{1}{n}\log\|D_{x}f^{n}(v)\|.$$
		\end{enumerate}
		
		Furthermore, if $f$ is a $C^{1}$ diffeomorphism, then  there exists an ~$f$ invariant set ~$Y$ $(i.e.~f(Y)=Y)$ with $\mu(Y)=1$ such that for each $x\in Y$,
		\begin{enumerate}
			\item There exists a $Df$ invariant splitting of the tangent space
			\begin{displaymath}
				T_{x}M=E_{1}(x)\oplus E_{2}(x)\oplus...\oplus E_{\ell(x)}(x).
			\end{displaymath}
			\item There exist measurable functions $\lambda_{1}(x)<\lambda_{2}(x)<...<\lambda_{\ell(x)}(x)$, such that for any ~$v\in E_{i}(x), ~v\neq 0$, we have~$$\lambda_{i}(x)=\lim_{n\rightarrow \infty}\frac{1}{n}\log\|D_{x}f^{n}(v)\|.$$
		\end{enumerate}
	\end{Thm}
	
	If $\mu$ is an ergodic measure for $f$, then $\lambda_{i}(x)$ and $\ell(x)$ are independent of the choice of the point $x$. Thus we can denote them by $\lambda_{1}(\mu), \lambda_{2}(\mu), \cdots,  \lambda_{\ell}(\mu)$ and we call these numbers the Lyapunov exponents of $f$ with respect to the measure $\mu$.
	
	%In this paper, for the notational convenience, we don't take the multiplicity of the Lyapunov exponents into account, that is, we denote the Lyapunov exponents of an ergodic measure $\mu$ as $\lambda_{1}(\mu)\leq\lambda_{2}(\mu)\leq...\leq\lambda_{d}(\mu)$.
	
	\fi
	\subsection{Sub-additive topological pressure and super-additive topological pressure.}
	
	Consider a continuous transformation $f: X\to X$ of a compact metric space $X$ equipped with a metric $d$. Denote by $\mathcal{M}(f)$ the set of all $f$-invariant Borel probability measures on $X$. We recall that a sequence of continuous functions (potentials) $\Phi=\{\varphi_n\}_{n\geq1}$ is called {\it sub-additive} if
	$$\varphi_{m+n}(x)\leq \varphi_n(x)+ \varphi_m\circ f^n(x),\ \  \forall m,n\in\mathbb{N} , \forall x \in X.$$
	
	
	Now we state the definition of sub-additive topological pressure which was introduced in \cite{cfh}.  For $x,y\in X$ and $n\geq 0$, define the metric $d_n$ on $X$ by
	$$d_n(x, y) = \max\{d(f^i(x),f^i(y))|0\leq i<n\}.$$
	Given $\varepsilon>0$ and $n\geq 0$, denote by $B_n(x,\varepsilon)= \{y\in X |d_n(x,y)<\varepsilon\}$ the Bowen's ball centered at $x$ of radius $\varepsilon$ and length $n$. We say a subset $E\subset X$ is {\it $(n,\varepsilon)$-separated} if $d_n(x,y)\geq\varepsilon$ for any two distinct points $x,y\in E$.
	
	Given a sub-additive sequence of continuous potentials $\Phi=\{\varphi_n\}_{n\geq 1}$, let
	$$P_n(\Phi,\varepsilon)= \sup\left\{\sum_{x\in E}e^{\varphi_n(x)}|E \mbox{ is an }(n,\varepsilon) \mbox{-separated set of }X\right\}.$$
	The quantity
	$$P(f,\Phi) =\lim_{\varepsilon\to 0}\limsup_{n\to\infty}\frac{1}{n}\log P_n(\Phi,\varepsilon)$$
	is called {\it the sub-additive topological pressure of $\Phi$}. One can show that it satisfies the following variational principle (see Theorem 1.1 of \cite{cfh}):
	\begin{equation}\label{vp}P(f,\Phi) =\sup_{\mu\in \mathcal{M}(f)}\left\{h_\mu(f)+\mathcal{F}_*(\Phi,\mu)|\mathcal{F}_*(\Phi,\mu)\neq\infty\right\},\end{equation}
	where $h_\mu(f)$ is the metric entropy of $f$ with respect to $\mu$ and
	$$\mathcal{F}_*(\Phi,\mu)=\lim_{n\to\infty}\frac{1}{n}\int\varphi_nd\mu.$$
	Existence of the above limit can be shown by the standard sub-additive argument.
	
	Next we recall the definition of the super-additive topological pressure introduced in \cite{caopesinzhao2019}.
	
	Given a sequence of super-additive continuous potentials $\Psi=\{\psi_{n}\}_{n\geq 1}$, that is a sequence of continuous functions satisfying
	\begin{align*}
		\psi_{m+n}(x)\geq \psi_{n}(f^{m}(x))+\psi_{m}(x)
	\end{align*}
	for any $n,m\in\mathbb{N}$ and any $x\in X$. The {\it super-additive topological pressure of $\Psi$} is defined as
	\begin{align*}
		P_{\text{var}}(f,\Psi)\triangleq \sup \{h_{\mu}(f)+\mathcal{F}_{\ast}(\Psi,\mu):\mu\in \mathcal{M}(f)\},
	\end{align*}
	where $\mathcal{F}_{\ast}(\Psi,\mu)=\lim_{n\rightarrow \infty}\frac{1}{n}\int \psi_{n}d\mu=\sup \frac{1}{n}\int \psi_{n}d\mu$. The second equality is due to the standard sub-additive argument.
	
	
	\subsection{Singular-valued potentials.}\label{sub}
	
	Suppose $M$ is a $d$-dimensional compact Riemannian manifold and $f:M\rightarrow M$ is a $C^{1+\gamma}$ expanding map with a repeller $\Lambda$ such that $f|_{\Lambda}$ is topological transitive. Given $x\in M$ and $n\geq 1$, we consider the differentiable operator $D_{x}f^{n}: T_{x}M\rightarrow T_{f^{n}(x)}M$ and denote its singular values in the decreasing order by
	$$\alpha_{1}(x,f^{n})\geq\alpha_{2}(x,f^{n})\geq\cdots\geq \alpha_{d}(x,f^{n}).$$
	For $s\in [0,d]$, set
	$$\varphi^{s}(x,f^{n})\triangleq\sum^{d}_{i=d-[s]+1} \log \alpha_{i}(x,f^{n})+(s-[s])\log\alpha_{d-[s]}(x,f^{n})$$
	and
	\begin{displaymath}
		\psi^{s}(x,f^{n})=\sum_{i=1}^{[s]}\log \alpha_{i}(x,f^{n})+(s-[s])\log \alpha_{[s]+1}(x,f^{n}),
	\end{displaymath}
	where $[s]$ is the largest integer not exceeding $s$. Since $f$ is a $C^{1+\gamma}$ map, the functions $x\mapsto \alpha_{i}(x,f^{n})$, $x\mapsto \varphi^{s}(x,f^{n})$ and $x\mapsto \psi^{s}(x,f^{n})$ are continuous.  It is easy to check that for all $m,n \in \mathbb{N}$,
	$$\varphi^{s}(x,f^{m+n})\geq \varphi^{s}(x,f^{n})+\varphi^{s}(f^n(x),f^{m}),$$
	\begin{align*}
		\psi^{s}(x,f^{m+n})\leq \psi^{s}(x,f^{n})+\psi^{s}(f^{n}(x),f^{m}).
	\end{align*}
	Thus $\Phi_{f}(s)\triangleq \{-\varphi^{s}(\cdot,f^{n})\}_{n\geq 1}$ are sub-additive and $\Psi_{f}(s)=\{-\psi^{s}(\cdot,f^{n})\}_{n\geq 0}$ are super-additive. We call them the {\it sub-additive singular-valued potentials} and the {\it super-additive singular-valued potentials }respectively.
	\iffalse
	The definitions of the singular valued potentials for hyperbolic diffeomorphisms are similar.
	
	Suppose $M$ is a $d$-dimensional compact Riemannian manifold and $f:M\rightarrow M$ is a $C^{1+\alpha}$ diffeomorphism admitting a hyperbolic set $\Lambda$ such that $f|_{\Lambda}$ is topological transitive. Given $x\in M$ and $n\geq 1$, we consider the differentiable operator $D_{x}f^{n}: T_{x}M\rightarrow T_{f^{n}(x)}M$ and denote its singular values in the decreasing order by
	$$\alpha_{1}(x,f^{n})\geq\alpha_{2}(x,f^{n})\geq\cdots\geq \alpha_{u}(x,f^{n})>1>\alpha_{u+1}(x,f^{n})\geq\cdots\geq \alpha_{d}(x,f^{n}),$$ where $u$ is the dimension of the unstable manifold.
	For $s\in [0,u]$, set
	$$\tilde{\varphi}^{s}(x,f^{n})\triangleq\sum^{u}_{i=u-[s]+1} \log \alpha_{i}(x,f^{n})+(s-[s])\log\alpha_{u-[s]}(x,f^{n}),$$
	and
	\begin{align*}
		\tilde{\psi}^{s}(x,f^{n})=\sum_{i=1}^{[s]}\log \alpha_{i}(x,f^{n})+(s-[s])\log \alpha_{[s]+1}(x,f^{n}).
	\end{align*}
	$\tilde{\Phi}_{f}(s)\triangleq \{-\tilde{\varphi}^{s}(\cdot,f^{n})\}_{n\geq 1}$ are sub-additive and $\tilde{\Psi}_{f}(s)=\{-\psi^{s}(\cdot,f^{n})\}_{n\geq 0}$ are super-additive. We call them the sub-additive singular valued potentials and the super-additive singular valued potentials respectively.
	\fi
	
	
	
	
	\subsection{Carath$\acute{\text{e}}$odory singular dimension.}\label{C dimension} We recall the definition of the Carath$\acute{\text{e}}$odory singular dimension defined in~\cite{caopesinzhao2019}, which is based on the Carath$\acute{\text{e}}$odory construction described in~\cite{Pesin}.
	
	Let $\Lambda$ be a repeller for an expanding map $f$ and $\Phi=\{-\varphi^{\alpha}(\cdot,f^{n})\}_{n\geq 1}$ be the sub-additive singular-valued potential defined in  subsection \ref{sub}. For every $Z\subset \Lambda$, define
	\begin{align*}
		m(Z,\alpha,r)\triangleq \lim_{N\rightarrow \infty} \inf\{\sum_{i}\exp(\sup_{y\in B_{n_{i}}(x_{i},r)} -\varphi^{\alpha}(y,f^{n_{i}}))\},
	\end{align*}
	where the infimum is taken over all the collections $\{B_{n_{i}}(x_{i},r)\}$ of Bowen's balls with $x_{i}\in \Lambda$, $n_{i}\geq N$ those cover $Z$. One can show that there is a jump-up value
	\begin{align*}
		\dim_{C,r} Z\triangleq \inf\{\alpha:m(Z,\alpha,r)=0\}=\sup\{\alpha:m(Z,\alpha,r)=+\infty\},
	\end{align*}
	which is called {\it the Carath$\acute{\text{e}}$odory singular dimension of $Z$}. For the motivation of this definition, one can refer to \cite{caopesinzhao2019}.
	\iffalse
	Similarly, one can define the Carath$\acute{\text{e}}$odory singular dimension of any subset of a hyperbolic set. For this, we only need to replace the sub-additive singular valued potential $\Phi$ by $\tilde{\Phi}$, which is defined in  subsection \ref{sub}.
	\fi
	
	\subsection{Hausdorff dimension}
	Now we recall the definition of the Hausdorff dimension, one can refer to \cite{Pesin}.
	
	Let $X$ be a compact metric space equipped with a metric $d$. Given  a subset $Z$ of $X$, for $s\geq 0$ and $\delta>0$, define
	\[
	\mathcal{H}_{\delta}^{s}(Z):=\inf \left\{\sum_{i}|U_i|^s: \
	Z\subset \bigcup_{i}U_i,~|U_i|\leq \delta\right\}
	\]
	where $|\cdot|$ denotes the diameter of a set. The quantity
	$\mathcal{H}^{s}(Z):=\lim\limits_{\delta\rightarrow 0}\mathcal{H}_{\delta}^{s}(Z)$ is called the {\em $s-$dimensional Hausdorff measure} of $Z$. Define the {\em Hausdorff dimension} of $Z$, denoted by $\dim_H  Z$, as follows:
	\[
	\dim_H  Z =\inf \{s:\ \mathcal{H}^{s}(Z)=0\}=\sup \{s: \mathcal{H}^{s}(Z)=\infty\}.
	\]
	One can check that the Hausdorff dimension satisfies the monotonicity, that is, for $Y_{1}\subset Y_{2}\subset X$, we have ~$\dim_{H}Y_{1}\leq \dim_{H}Y_{2}$.
	
	
	Given a probability Borel measure $\mu$ on $X$, the {\it Hausdorff dimension of the measure $\mu$}  is defined as $$\dim_{H}\mu=\inf\{\dim_{H}Y:Y\subset X ,  \mu(Y)=1\}.$$ This quantity was introduced by Young in \cite{Young}.
	\subsection{Lower bound for the Hausdorff dimension of a non-conformal repeller.}
	Recently, Cao, Pesin and Zhao~\cite{caopesinzhao2019} gave a lower bound for the Hausdorff dimension of a non-conformal repeller, which is given by the root of the Bowen's equation involving the super-additive topological pressure of the super-additive singular valued potential. Take $\Psi_{f}(s)=\{-\psi^{s}(\cdot,f^{n})\}_{n\geq 0}$ as the super-additive singular-valued potentials and  $P_{\text{sup}}(s)=P_{\text{var}}(f|_{\Lambda},\Psi_{f}(s))$ as the super-additive topological pressure of the super-additive singular-valued potential.
	\begin{Thm}\cite{caopesinzhao2019}\label{lower bound}
		Let $\Lambda$ be a repeller for a $C^{1+\gamma}$ expanding map $f:M\rightarrow M$. Then
		\begin{align*}
			dim_{H}\Lambda\geq s^{\ast},
		\end{align*}
		where $s^{\ast}$ is the unique root of the equation $P_{\text{sup}}(s)=0$.
	\end{Thm}
	
	
	
	\subsection{The Carath$\acute{\text{e}}$odory singular dimension of a repeller}
	In~\cite{caopesinzhao2019}, Cao, Pesin and Zhao obtained that the Carath$\acute{\text{e}}$odory singular dimension of a repeller for a $C^{1+\gamma}$ map is exactly the unique root of the Bowen's equation involving the sub-additive topological pressure of the sub-additive singular valued potential. Take $\Phi_{f}(\alpha)=\{-\varphi^{\alpha}(\cdot,f^{n})\}_{n\geq 0}$ as the sub-additive singular valued potentials and denote by $P(f|_{\Lambda},\alpha)=P(f|_{\Lambda},\Phi_{f}(\alpha))$ the sub-additive topological pressure of the sub-additive singular-valued potential.
	\begin{Thm}\cite{caopesinzhao2019}\label{singular-valued dimension}
		Let $f: M\rightarrow M$ be a $C^{1+\gamma}$ expanding map with a repeller $\Lambda$. Assume that $f|_{\Lambda}$ is topologically transitive. Then $dim_{C,r} \Lambda=\alpha_{0}$ for all sufficiently small $r>0$, where $\alpha_{0}$ is the unique root of the Bowen's equation $P(f|_{\Lambda},\alpha)=0$.
	\end{Thm}
	\begin{Rem}
		As a result, we can use the notation $dim_{C}\Lambda$ to express the Carath$\acute{\text{e}}$odory singular dimension of the repeller $\Lambda$.
	\end{Rem}
	
	
	\subsection{Approximate of an hyperbolic measure by repellers} Results of the following type is usually referred to Katok~\cite{KH}.
	\begin{Thm}\label{approximate}
		Suppose $M$ is a $d-$dimensional Riemannian manifold. Let $f: M\rightarrow M$ be a $C^{1+\gamma}$ expanding map with a repeller $\Lambda$ and $\mu$ be an ergodic measure on $\Lambda$ with $h_{\mu}(f)>0$. Then for any $\varepsilon>0$, there exists an $f-$invariant compact subset $\Lambda_{\varepsilon}\subset\Lambda$ such that $h_{top}(f|_{\Lambda_{\varepsilon}})\geq h_{\mu}(f)-\varepsilon$ and for any $f-$invariant ergodic measure $\nu$ supported on $\Lambda_{\varepsilon}$, it holds that $|\lambda_{i}(\mu)-\lambda_{i}(\nu)|\leq\varepsilon$ for each $1\leq i\leq d$.
	\end{Thm}
	
	\iffalse
	\begin{Thm}\label{hyperbolic approximate}
		Suppose $M$ is a $d-$dimensional Riemannian manifold. Let $f: M\rightarrow M$ be a $C^{1+\alpha}$ diffeomorphism admitting a hyperbolic set $\Lambda$ and $\mu$ be an ergodic measure on $\Lambda$ with $h_{\mu}(f)>0$. Then for any $\varepsilon>0$, there exists an $f-$invariant compact subset $\Lambda_{\varepsilon}\subset\Lambda$ such that $h_{top}(f|_{\Lambda_{\varepsilon}})\geq h_{\mu}(f)-\varepsilon$ and for any $f-$invariant measure $\nu$ supported on $\Lambda_{\varepsilon}$, it holds that $|\lambda_{i}(\mu)-\lambda_{i}(\mu)|\leq\varepsilon$ for each $1\leq i\leq d$.
	\end{Thm}
	
	\fi
	
	
	
	
	
	
	
	
	\iffalse
	\section{Dimension of the unstable manifold}
	In this section, we establish the corresponding results of Theorem \ref{lower bound} and Theorem \ref{singular-valued dimension} for the dimension of the intersection of the hyperbolic set and the unstable manifold.
	
	First we recall the definition of the dominated splitting.
	
	Let $M$ be a $d-$dimensional compact smooth Riemannian manifold and $f:M\rightarrow M$ be a $C^{1+\alpha}$ diffeomorphism. Suppose $\Lambda\subset M$ is a compact $f-$invariant set. We say $\Lambda$ admits a dominated splitting if there is a continuous invariant splitting $T_{\Lambda}M=E\oplus F$ and constants $C>0,\lambda\in (0,1)$ such that for each $x\in \Lambda$, $n\in \mathbb{N}$, $0\neq u\in E(x)$ and $0\neq v\in F(x)$
	\begin{align*}
		\frac{\|D_{x}f^{n}(u)\|}{\|u\|}\leq C\lambda^{n}\frac{\|D_{x}f^{n}(v)\|}{\|v\|}.
	\end{align*}
	We write $E\preceq F$ if $F$ dominates E. Furthermore, given $0<\ell\leq d$, we say a continuous invariant splitting $T_{\Lambda}M=E_{1}\oplus...\oplus E_{\ell}$ dominated if there are numbers $\lambda_{1}<\lambda_{2}<...<\lambda_{\ell}$, constants $C>0$ and $0<\varepsilon<\min_{i}\frac{\lambda_{i+1}-\lambda_{i}}{100}$ such that for every $x\in \Lambda$, $n\in\mathbb{N}$ and $1\leq j\leq \ell$ and each unit vector $u\in E_{j}(x)$
	\begin{align*}
		C^{-1}\text{exp}(n(\lambda_{j}-\varepsilon))\leq \|D_{x}f^{n}(u)\|\leq C \text{exp}(n(\lambda_{j}+\varepsilon)).
	\end{align*}
	In particular, $E_{1}\preceq...\preceq E_{\ell}$. We shall use the notion $\{\lambda_{j}\}-$dominated when we want to stress the dependence on the numbers $\{\lambda_{j}\}$.
	
	The authors in \cite{ACW} gave the proof of the following result, which enables us to give a sharper lower bound of the Hausdorff dimension of the intersection of the hyperbolic set and the unstable manifold.
	\begin{Thm}\cite{ACW}\label{Katok}
		Let $f$ be a $C^{1+\alpha}$ diffeomorphism on a compact Riemannian manifold $M$ and $\Lambda\subset M$ is a locally maximal hyperbolic set. Suppose $\mu$ is an ergodic invariant measure on $\Lambda$ with $h_{\mu}(f)>0$. Then for any $\varepsilon>0$, there exists an $f-$invariant compact subset $\Lambda_{\varepsilon}\subset \Lambda$ such that
		\begin{enumerate}
			\item $h_{\text{top}}(f|_{\Lambda_{\varepsilon}})\geq h_{\mu}(f)-\varepsilon$,
			\item there is  a $\{\lambda_{j}(\mu)\}$-dominated splitting $T_{x}M=E_{1}\oplus E_{2}\oplus...\oplus E_{\ell}$ over $\Lambda_{\varepsilon}$, where $\lambda_{1}(\mu)<...<\lambda_{\ell}(\mu)$ are distinct Lyapunov exponents of $f$ with respect to the measure $\mu$.
		\end{enumerate}
	\end{Thm}
	
	
	
	
	Assume that the map $f|_{\Lambda}$ possesses a $\{\lambda_{j}\}$-dominated splitting $T_{\Lambda}M=E_{1}\oplus E_{2}\oplus...\oplus E_{k}$ with $E_{1}\succeq E_{2}\succeq ...\succeq E_{k}$ and $\lambda_{1}>\lambda_{2}>...>\lambda_{\ell}>0>\lambda_{\ell+1}>...>\lambda_{k}$. Let $m_{j}=dim E_{j}$, $r_{j}=m_{1}+...+m_{j}$ for $j\in\{1,2,...,\ell\}$ and $r_{0}=0$. Set $u=m_{1}+...+m_{\ell}$. For $s\in [0,u]$, $n\geq 1$ and $x\in \Lambda$, define
	\begin{align*}
		\hat{\psi}^{s}(x,f^{n})=\sum_{j=1}^{d}m_{j}\log\|D_{x}f^{n}|_{E_{j}}\|+(s-r_{d})\log\|D_{x}f^{n}|_{E_{d+1}}\|
	\end{align*}
	if $r_{d}\leq s<r_{d+1}$ for some $d\in \{0,1,...,\ell-1\}$. It is clear that $\hat{\Psi}_{f}(s)\triangleq \{-\hat{\psi}^{s}(x,f^{n})\}_{n\geq 1}$ is super-additive.
	%Denote $\hat{P}_{sup}(s)\triangleq P_{var}(f|_{\Lambda},\hat{\Psi}_{f}(s))$. It is easy to see that $\hat{P}_{sup}(s)$ is continuous and strictly decreasing in $s$.
	
	
	\begin{Lem}\label{lower bound 1}
		Assume that the map $f|_{\Lambda}$ possesses a $\{\lambda_{j}\}$-dominated splitting $T_{\Lambda}M=E_{1}\oplus E_{2}\oplus...\oplus E_{k}$ with $E_{1}\succeq E_{2}\succeq ...\succeq E_{k}$ and $\lambda_{1}>\lambda_{2}>...>\lambda_{\ell}>0>\lambda_{\ell+1}>...>\lambda_{k}$. Then for any $x\in \Lambda$, we have $dim_{H}\Lambda\cap W^{u}_{\text{loc}}(x)\geq s_{1}$, where $s_{1}$ is the unique root of Bowen's equation
		\begin{align*}
			P(f|_{\Lambda},-\hat{\psi}^{s}(\cdot,f))=0.
		\end{align*}
	\end{Lem}
	
	
	\begin{proof}
		Let $\{P_{1},P_{2},...,P_{l}\}$ be a Markov partition of $\Lambda$. Since $\Lambda$ is locally maximal, there exists $\delta>0$ such that for each $i=1,2,...,l$, the closed $\delta-$neighborhood $\tilde{P}_{i}$ of $P_{i}$ satisfies $\tilde{P}_{i}\subset U$, where $U$ is an open neighborhood of $\Lambda$  in the definition of the locally maximal hyperbolic set, and $\tilde{P}_{i}\cap \tilde{P}_{j}=\varnothing$ whenever $P_{i}\cap P_{j}=\varnothing$. Given an admissible sequence $\mathbf{i}=...i_{-m}...i_{0}i_{1}...i_{n-1}$ and the cylinder $P_{...i_{-m}...i_{0}i_{1}...i_{n-1}}$, we denote by $\tilde{P}_{...i_{-m}...i_{0}i_{1}...i_{n-1}}$ the corresponding cylinder. Note that the $\{\lambda_{j}\}$-dominated splitting can be extended to $U$ and the map $x\mapsto E_{i}(x)$ is H$\ddot{\text{o}}$lder continuous on $\Lambda$ and also can be extended to $U$. So is the map $x\mapsto \|D_{x}f|_{E_{i}}\|$ for each $i=1,2,...,k.$
		
		Since $\hat{\psi}^{s_{1}}(\cdot,f)$ is H$\ddot{\text{o}}$lder continuous on $\Lambda$, there is a unique equilibrium state $\mu$ for $P(f|_{\Lambda},-\hat{\psi}^{s_{1}}(\cdot,f))$, which is a Gibbs measure. Fix $x\in\Lambda$, the conditional measure
		$\mu_{x}^{u}$ has the $u-$Gibbs property. Therefore, if $\pi(...\bar{i}_{-m}...\bar{i}_{0}...\bar{i}_{n}...)=x$, then there exists $K>0$ such that for each $n\in \mathbb{N}$ and each $y\in P_{...\bar{i}_{-m}...\bar{i}_{0}\bar{i}_{1}...\bar{i}_{n-1}}$, we have
		\begin{align*}
			K^{-1}\exp(-\sum_{i=0}^{n-1}\hat{\psi}^{s_{1}}(f^{i}(y),f))\leq \mu_{x}^{u}(P_{...\bar{i}_{-m}...\bar{i}_{0}\bar{i}_{1}...\bar{i}_{n-1}})\leq K \exp(-\sum_{i=0}^{n-1}\hat{\psi}^{s_{1}}(f^{i}(y),f)).
		\end{align*}
		Since $T_{\Lambda}M=E_{1}\oplus E_{2}\oplus...\oplus E_{k}$ is a dominated splitting, the angles between different subspaces $E_{i}$ are uniformly bounded away from zero. Therefore there exists $a>0$ such that for each admissible sequence $...i_{-m}...i_{0}i_{1}...i_{n-1}$, there exists $\xi_{i}\in \tilde{P}_{...i_{-m}...i_{0}i_{1}...i_{n-1}}$ for each $i=1,2,...,\ell$ with a rectangle of sides
		\begin{align*}
			\overbrace{a\|D_{\xi_{1}}f^{n}|_{E_{1}}\|^{-1},...,a\|D_{\xi_{1}}f^{n}|_{E_{1}}\|^{-1},}^{m_{1}}...,\overbrace{a\|D_{\xi_{\ell}}f^{n}|_{E_{\ell}}\|^{-1},...,a\|D_{\xi_{\ell}}f^{n}|_{E_{\ell}}\|^{-1}}^{m_{\ell}}
		\end{align*}
		contained in $\tilde{P}_{...i_{-m}...i_{0}i_{1}...i_{n-1}}$. One can show that there exists $C_{0}>0$ such that for any $\tilde{P}_{...i_{-m}...i_{0}i_{1}...i_{n-1}}$ and any $y, z \in \tilde{P}_{...i_{-m}...i_{0}i_{1}...i_{n-1}}$, any $i\in \{1,2,...,\ell\}$, we have
		\begin{align*}
			\frac{1}{C_{0}}\leq \frac{\prod_{j=0}^{n-1}\|D_{f^{j}(y)}f|_{E_{i}}\|^{-1}}{\prod_{j=0}^{n-1}\|D_{f^{j}(z)}f|_{E_{i}}\|^{-1}}\leq C_{0}.
		\end{align*}
		Notice that $\|D_{\xi_{i}}f^{n}|_{E_{i}}\|^{-1}\geq \prod_{j=0}^{n-1}\|D_{f^{j}(\xi_{i})}f|_{E_{i}}\|^{-1}$. Thus there exists $a_{1}>0$ and $y\in P_{...i_{-m}...i_{0}i_{1}...i_{n-1}}$ such that $\tilde{P}_{...i_{-m}...i_{0}i_{1}...i_{n-1}}$ contains a rectangle of sides
		\begin{align*}
			\overbrace{a_{1}A_{1}(y,n),...,a_{1}A_{1}(y,n),}^{m_{1}}\overbrace{a_{1}A_{2}(y,n),...,a_{1}A_{2}(y,n),}^{m_{2}}...,\\
			\overbrace{a_{1}A_{\ell}(y,n),...,a_{1}A_{\ell}(y,n)}^{m_{\ell}},
		\end{align*}
		where $A_{i}(y,n)=\prod_{j=0}^{n-1}\|D_{f^{j}(y)}f|_{E_{i}}\|^{-1}$ for each $i\in \{1,2,...,\ell\}$.
		Notice that there is $i\in \{1,2,...,\ell\}$ such that $r_{i}\leq s_{1}<r_{i+1}$. Fix $r>0$ small enough and set
		\begin{align*}
			\mathcal{Q}=\{\mathbf{i}=(...\bar{i}_{-m}...\bar{i}_{0}i_{1}...i_{n-1}):a_{1}A_{i+1}(y,n)\leq r\ \text{for\ all}\ y \in P_{...\bar{i}_{-m}...\bar{i}_{0}i_{1}...i_{n-1}};\\
			\text{but}~a_{1}A_{i+1}(z,n-1)>r~\text{for~some}~z \in P_{...\bar{i}_{-m}...\bar{i}_{0}i_{1}...i_{n-1}}\}.
		\end{align*}
		Therefore, for each $\mathbf{i}=(...\bar{i}_{-m}...\bar{i}_{0}i_{1}...i_{n-1})\in \mathcal{Q}$, we have
		\begin{align*}
			br<a_{1}A_{i+1}(y,n)\leq r ~\text{for ~all}~y\in P_{...\bar{i}_{-m}...\bar{i}_{0}i_{1}...i_{n-1}},
		\end{align*}
		where $b=C_{0}^{-1}\min_{x\in \Lambda}\|D_{x}f\|^{-1}$.
		
		Let $B\subset W^{u}_{\text{loc}}(x)$ be a ball of radius $r$ and $\tilde{B}\subset W^{u}_{\text{loc}}(x)$ a ball of radius $Cr$ with the same center as that of $B$, where $C$ is a constant only depended on the dimension of the unstable manifold. Put
		\begin{align*}
			\mathcal{Q}_{1}=\{\mathbf{i}\in \mathcal{Q}|P_{\mathbf{i}}\cap B\neq \varnothing\}.
		\end{align*}
		Hence, for each $\mathbf{i}\in\mathcal{Q}_{1}$, we have $\tilde{P}_{\mathbf{i}}\cap \tilde{B}$ contains a rectangle of sides
		\begin{align*}
			\overbrace{a_{1}A_{1}(y,n),...,a_{1}A_{1}(y,n)}^{m_{1}},...,\overbrace{a_{1}A_{i}(y,n),...,a_{1}A_{i}(y,n)}^{m_{i}},\\
			\overbrace{a_{1}A_{i+1}(y,n),...,a_{1}A_{i+1}(y,n)}^{u-r_{i}}.
		\end{align*}
		It follows that
		\begin{align*}
			a_{1}^{u}A_{i+1}(y,n)^{u-r_{i}}A_{i}(y,n)^{m_{i}}...A_{1}(y,n)^{m_{1}}\leq vol (\tilde{P}_{\mathbf{i}}\cap \tilde{B}).
		\end{align*}
		Since
		\begin{align*}
			A_{i+1}(y,n)^{u-r_{i}}&=A_{i+1}(y,n)^{s_{1}-r_{i}}A_{i+1}(y,n)^{u-s_{1}}\\
			&\geq A_{i+1}(y,n)^{s_{1}-r_{i}}(\frac{b}{a_{1}})^{u-s_{1}}r^{u-s_{1}},
		\end{align*}
		we have
		\begin{align*}
			a_{2}r^{u-s_{1}}A_{i+1}(y,n)^{s_{1}-r_{i}}A_{i}(y,n)^{m_{i}}...A_{1}(y,n)^{m_{1}}\leq vol (\tilde{P}_{\mathbf{i}}\cap \tilde{B})
		\end{align*}
		for some constant $a_{2}>0$. Therefore,
		\begin{align*}
			\sum_{\mathbf{i}\in\mathcal{Q}_{1}}a_{2}r^{u-s_{1}}A_{i+1}(y,n)^{s_{1}-r_{i}}A_{i}(y,n)^{m_{i}}...A_{1}(y,n)^{m_{1}}\leq vol (\tilde{B})\leq C^{u}a_{3}r^{u}
		\end{align*}
		for some constant $a_{3}>0$. Hence,
		\begin{align*}
			\sum_{\mathbf{i}\in\mathcal{Q}_{1}}A_{i+1}(y,n)^{s_{1}-r_{i}}A_{i}(y,n)^{m_{i}}...A_{1}(y,n)^{m_{1}}\leq a_{4}r^{s_{1}}
		\end{align*}
		for some constant $a_{4}>0$.
		On the other hand, we have
		\begin{align*}
			\mu_{x}^{u}(B)&\leq \sum_{(...\bar{i}_{-m}...\bar{i}_{0}i_{1}...i_{n-1})\in\mathcal{Q}_{1}}\mu_{x}^{u}(P_{...\bar{i}_{-m}...\bar{i}_{0}i_{1}...i_{n-1}})\\
			&\leq K \sum_{(...\bar{i}_{-m}...\bar{i}_{0}i_{1}...i_{n-1})\in\mathcal{Q}_{1}} exp(-\sum_{i=0}^{n-1}\hat{\psi}^{s_{1}}(f^{i}(y),f))\\
			&=K \sum_{(...\bar{i}_{-m}...\bar{i}_{0}i_{1}...i_{n-1})\in\mathcal{Q}_{1}} A_{i+1}(y,n)^{s_{1}-r_{i}}A_{i}(y,n)^{m_{i}}...A_{1}(y,n)^{m_{1}}\\
			&\leq a_{5}r^{s_{1}}
		\end{align*}
		for some constant $a_{5}>0$. This implies that $\dim_{H}\mu_{x}^{u}\geq s_{1}$ and hence $\dim_{H}\Lambda\cap W^{u}_{\text{loc}}(x)\geq \dim_{H}\mu_{x}^{u}\geq s_{1}$.
	\end{proof}
	
	
	\iffalse
	\begin{proof}
		Let $\{P_{1},P_{2},...,P_{l}\}$ be a Markov partition of $\Lambda$. Since $\Lambda$ is locally maximal, there exists $\delta>0$ such that for each $i=1,2,...,l$, the closed $\delta-$neighborhood $\tilde{P}_{i}$ of $P_{i}$ satisfies $\tilde{P}_{i}\subset U$, where $U$ is an open neighborhood of $\Lambda$  in the definition of the locally maximal hyperbolic set, and $\tilde{P}_{i}\cap \tilde{P}_{j}=\varnothing$ whenever $P_{i}\cap P_{j}=\varnothing$. Given an admissible sequence $\mathbf{i}=i_{0}i_{1}...i_{n-1}$ and the cylinder $P_{i_{0}i_{1}...i_{n-1}}$, we denote by $\tilde{P}_{i_{0}i_{1}...i_{n-1}}$ the corresponding cylinder. Note that the $\{\lambda_{j}\}$-dominated splitting can be extended to $U$ and the map $x\mapsto E_{i}(x)$ is H$\ddot{o}$lder continuous on $\Lambda$ and also can be extended to $U$. So is the map $x\mapsto \|D_{x}f|_{E_{i}}\|$ for each $i=1,2,...,k.$
		
		Since $\hat{\psi}^{s_{1}}(\cdot,f)$ is H$\ddot{o}$lder continuous on $\Lambda$, there is a unique equilibrium state $\mu$ for $P(f|_{\Lambda},-\hat{\psi}^{s_{1}}(\cdot,f))$, which is a Gibbs measure. For $x\in\Lambda$, the conditional measure
		$\mu_{x}^{u}$ has the $u-$Gibbs property. Therefore, if $\pi(...i_{-n}...i_{0}...i_{n}...)=x$, then there exists $K>0$ such that for each $n\in \mathbb{N}$ and each $y\in P_{i_{0}i_{1}...i_{n-1}}$, we have
		\begin{align*}
			K^{-1}exp(-\sum_{i=0}^{n-1}\hat{\psi}^{s_{1}}(f^{i}(y),f))\leq \mu_{x}^{u}(P_{i_{0}i_{1}...i_{n-1}})\leq K exp(-\sum_{i=0}^{n-1}\hat{\psi}^{s_{1}}(f^{i}(y),f)).
		\end{align*}
		Since $T_{\Lambda}M=E_{1}\oplus E_{2}\oplus...\oplus E_{k}$ is a dominated splitting, the angles between different subspaces $E_{i}$ are uniformly bounded away from zero. Therefore there exists $a>0$ such that for each admissible sequence $i_{0}i_{1}...i_{n-1}$, there exists $\xi_{i}\in \tilde{P}_{i_{0}i_{1}...i_{n-1}}$ for each $i=1,2,...,\ell$ with a rectangle of sides
		\begin{align*}
			\overbrace{a\|D_{\xi_{1}}f^{n}|_{E_{1}}\|^{-1},...,a\|D_{\xi_{1}}f^{n}|_{E_{1}}\|^{-1},}^{m_{1}}...,\overbrace{a\|D_{\xi_{k}}f^{n}|_{E_{k}}\|^{-1},...,a\|D_{\xi_{k}}f^{n}|_{E_{k}}\|^{-1}}^{m_{\ell}},\overbrace{a,a,...,a}^{d-u}
		\end{align*}
		contained in $\tilde{P}_{i_{0}i_{1}...i_{n-1}}$. One can show that there exists $C_{0}>0$ such that for any $\tilde{P}_{i_{0}i_{1}...i_{n-1}}$ and any $y, z \in \tilde{P}_{i_{0}i_{1}...i_{n-1}}$, any $i\in \{1,2,...,\ell\}$, we have
		\begin{align*}
			\frac{1}{C_{0}}\leq \frac{\prod_{j=0}^{n-1}\|D_{f^{j}(y)}f|_{E_{i}}\|^{-1}}{\prod_{j=0}^{n-1}\|D_{f^{j}(z)}f|_{E_{i}}\|^{-1}}\leq C_{0}.
		\end{align*}
		Notice that $\|D_{\xi_{i}}f^{n}|_{E_{i}}\|^{-1}\geq \prod_{j=0}^{n-1}\|D_{f^{j}(\xi_{i})}f|_{E_{i}}\|^{-1}$. Thus there exists $a_{1}>0$ and $y\in P_{i_{0}i_{1}...i_{n-1}}$ such that $\tilde{P}_{i_{0}i_{1}...i_{n-1}}$ contains a rectangle of sides
		\begin{align*}
			\overbrace{a_{1}A_{1}(y,n),...,a_{1}A_{1}(y,n),}^{m_{1}}\overbrace{a_{1}A_{2}(y,n),...,a_{1}A_{2}(y,n),}^{m_{2}}...,\\
			\overbrace{a_{1}A_{\ell}(y,n),...,a_{1}A_{\ell}(y,n)}^{m_{\ell}},\overbrace{a_{1},a_{1},...,a_{1}}^{d-u},
		\end{align*}
		where $A_{i}(y,n)=\prod_{j=0}^{n-1}\|D_{f^{j}(y)}f|_{E_{i}}\|^{-1}$ for each $i\in \{1,2,...,\ell\}$.
		Notice that there is $i\in \{1,2,...,\ell\}$ such that $r_{i}\leq s_{1}<r_{i+1}$. Fix $r>0$ small enough and set
		\begin{align*}
			\mathcal{Q}=\{\mathbf{i}=(i_{0}i_{1}...i_{n-1}):a_{1}A_{i+1}(y,n)\leq r\ for\ all\ y \in P_{i_{0}i_{1}...i_{n-1}};\\
			but~a_{1}A_{i+1}(z,n-1)>r~for~some~z \in P_{i_{0}i_{1}...i_{n-1}}\}.
		\end{align*}
		Therefore, for each $\mathbf{i}=(i_{0}i_{1}...i_{n-1})\in \mathcal{Q}$, we have
		\begin{align*}
			br<a_{1}A_{i+1}(y,n)\leq r ~for ~all~y\in P_{i_{0}i_{1}...i_{n-1}},
		\end{align*}
		where $b=C_{0}^{-1}min_{x\in \Lambda}\|D_{x}f\|^{-1}$.
		
		Let $B$ be a ball of radius $r$ and $\tilde{B}$ a ball of radius $2r$. Put
		\begin{align*}
			\mathcal{Q}_{1}=\{\mathbf{i}\in \mathcal{Q}|P_{\mathbf{i}}\cap B\neq \varnothing\}.
		\end{align*}
		Hence, for each $\mathbf{i}\in\mathcal{Q}_{1}$, we have $\tilde{P}_{\mathbf{i}}\cap \tilde{B}$ contains a rectangle of sides
		\begin{align*}
			\overbrace{a_{1}A_{1}(y,n),...,a_{1}A_{1}(y,n)}^{m_{1}},...,\overbrace{a_{1}A_{i}(y,n),...,a_{1}A_{i}(y,n)}^{m_{i}},\\
			\overbrace{a_{1}A_{i+1}(y,n),...,a_{1}A_{i+1}(y,n)}^{u-r_{i}},\overbrace{a_{1},a_{1},...,a_{1}}^{d-u}.
		\end{align*}
		It follows that
		\begin{align*}
			a_{1}^{d}A_{i+1}(y,n)^{u-r_{i}}A_{i}(y,n)^{m_{i}}...A_{1}(y,n)^{m_{1}}\leq vol (\tilde{P}_{\mathbf{i}}\cap \tilde{B}).
		\end{align*}
		Since
		\begin{align*}
			A_{i+1}(y,n)^{u-r_{i}}&=A_{i+1}(y,n)^{s_{1}-r_{i}}A_{i+1}(y,n)^{u-s_{1}}\\
			&\geq A_{i+1}(y,n)^{s_{1}-r_{i}}(\frac{b}{a_{1}})^{u-s_{1}}r^{u-s_{1}},
		\end{align*}
		we have
		\begin{align*}
			a_{2}r^{u-s_{1}}A_{i+1}(y,n)^{s_{1}-r_{i}}A_{i}(y,n)^{m_{i}}...A_{1}(y,n)^{m_{1}}\leq vol (\tilde{P}_{\mathbf{i}}\cap \tilde{B})
		\end{align*}
		for some constant $a_{2}>0$. Therefore,
		\begin{align*}
			\sum_{\mathbf{i}\in\mathcal{Q}_{1}}a_{2}r^{u-s_{1}}A_{i+1}(y,n)^{s_{1}-r_{i}}A_{i}(y,n)^{m_{i}}...A_{1}(y,n)^{m_{1}}\leq vol (\tilde{B})\leq 2^{d}a_{3}r^{d}
		\end{align*}
		for some constant $a_{3}>0$. Hence,
		\begin{align*}
			\sum_{\mathbf{i}\in\mathcal{Q}_{1}}A_{i+1}(y,n)^{s_{1}-r_{i}}A_{i}(y,n)^{m_{i}}...A_{1}(y,n)^{m_{1}}\leq a_{4}r^{s_{1}}
		\end{align*}
		for some constant $a_{4}>0$.
		On the other hand, we have
		\begin{align*}
			\mu_{x}(B)&\leq \sum_{(i_{0}i_{1}...i_{n-1})\in\mathcal{Q}_{1}}\mu_{x}(P_{i_{0}i_{1}...i_{n-1}})\\
			&\leq K \sum_{(i_{0}i_{1}...i_{n-1})\in\mathcal{Q}_{1}} exp(-\sum_{i=0}^{n-1}\hat{\psi}^{s_{1}}(f^{i}(y),f))\\
			&=K \sum_{(i_{0}i_{1}...i_{n-1})\in\mathcal{Q}_{1}} A_{i+1}(y,n)^{s_{1}-r_{i}}A_{i}(y,n)^{m_{i}}...A_{1}(y,n)^{m_{1}}\\
			&\leq a_{5}r^{s_{1}}
		\end{align*}
		for some constant $a_{5}>0$. This implies that $dim_{H}\mu_{x}\geq s_{1}$ and hence $dim_{H}\Lambda\cap W^{u}_{\text{loc}}(x,f)\geq dim_{H}\mu_{x}\geq s_{1}$.
	\end{proof}
	
	\fi
	
	
	With the help of Theorem \ref{Katok} and Lemma \ref{lower bound 1}, one can follow the proof in \cite{caopesinzhao2019} step by step to prove the following result, so we omit the details.
	\begin{Thm}\label{lower bound for hyperbolic}
		Let $f: M\rightarrow M$ be a $C^{1+\alpha}$ diffeomorphism with a hyperbolic set $\Lambda$. Assume that $f|_{\Lambda}$ is topologically transitive.  Then for any $x\in \Lambda$, we have
		\begin{align*}
			dim_{H}\Lambda\cap W^{u}_{\text{loc}}(x)\geq s_{\ast},
		\end{align*}
		where $s_{\ast}$ is the unique root of the equation $P_{var}(f|_{\Lambda},\tilde{\Psi}_{f}(s))=0$.
	\end{Thm}
	
	The proof of the following result is the same as the proof of the result for repellers in \cite{caopesinzhao2019}, so we omit it.
	\begin{Thm}\label{singular-valued dimension for hyperbolic}
		Let $f: M\rightarrow M$ be a $C^{1+\alpha}$ diffeomorphism with a hyperbolic set $\Lambda$. Assume that $f|_{\Lambda}$ is topologically transitive. Then for any $x\in \Lambda$, we have $$dim_{C,r} \Lambda \cap W^{u}_{\text{loc}}(x)=\alpha_{0}$$ for all sufficiently small $r>0$, where $\alpha_{0}$ is the unique root of the Bowen's equation $P(f|_{\Lambda},\tilde{\Phi}_{f}(\alpha))=0$.
	\end{Thm}
	
	
	\begin{proof}
		First we prove that for any $0<t<\alpha_{0}$, we have $\dim_{C}\Lambda\cap W^{u}(x)> t$.
		
		Since $P(f|_{\Lambda},\tilde{\Phi}_{f}(t))$ is strictly decreasing in $t$, fix $0<t<\alpha_{0}$, we have  $P(f|_{\Lambda},\tilde{\Phi}_{f}(t))>0$. And we take $\delta>0$ such that $P(f|_{\Lambda},\tilde{\Phi}_{f}(t))-\delta>0$. Since $P(f|_{\Lambda},\tilde{\Phi}_{f}(t))=\lim_{r\rightarrow 0}P(f|_{\Lambda},\tilde{\Phi}_{f}(t),r)$, there exists $r_{0}>0$ such that for each $0<r<r_{0}$ one has
		$$P(f|_{\Lambda},\tilde{\Phi}_{f}(t))-\delta<P(f|_{\Lambda},\tilde{\Phi}_{f}(t),r)<P(f|_{\Lambda},\tilde{\Phi}_{f}(t))+\delta.$$
		Fix such a number $r>0$, by the hyperbolicity and transitivity of $\Lambda$, there exists $m\in\mathbb{N}$, such that for any $y\in\Lambda$, we have
		$$W_{\frac{r}{2}}^{s}(y)\cap f^{m}W_{r}^{u}(x)\neq \varnothing.$$
		
		
		Fix $C\geq 1$ and $0<\chi<1$ as the skewness in the definition of the hyperbolicity. And take $k\in \mathbb{N}$ such that $C\chi^{k}\leq \frac{1}{2}$.
		Suppose $\{B_{n_{i}}(x_{i},r)\}$ is a cover of $\Lambda\cap W^{u}_{r}(x)$ consisting of Bowen's balls with $x_{i}\in \Lambda\cap W^{u}_{r}(x)$, then $\{f^{m}B_{n_{i}}(x_{i},r)\}$ cover $\Lambda\cap f^{m}W^{u}_{r}(x)$. Notice that $f^{m}B_{n_{i}}(x_{i},r)\subset B_{n_{i}-m}(f^{m}(x_{i}),r)$. Thus $\{B_{n_{i}-m}(f^{m}(x_{i}),r)\}$ is a cover of $\Lambda\cap f^{m}W^{u}_{r}(x)$ consisting of Bowen's balls with $f^{m}(x_{i})\in \Lambda\cap f^{m}W^{u}_{r}(x)$. For any $y\in\Lambda$, recall that $$W_{\frac{r}{2}}^{s}(y)\cap f^{m}W_{r}^{u}(x)\neq \varnothing.$$
		Thus there exists $z\in W_{\frac{r}{2}}^{s}(y)\cap f^{m}W_{r}^{u}(x)\cap B_{n_{i}-m}(f^{m}(x_{i}),r)$ for some $B_{n_{i}-m}(f^{m}(x_{i}),r)\in \{B_{n_{i}-m}(f^{m}(x_{i}),r)\}$. And hence
		\begin{align*}
			d(y,f^{m}(x_{i}))\leq d(y,z)+d(z,f^{m}(x_{i}))\leq r.
		\end{align*}
		Therefore, $\{B_{n_{i}-m}(f^{m}(x_{i}),r)\}$ is a cover of $\Lambda$.
		
		By the sub-additivity of $\{-\tilde{\varphi}^{t}(\cdot,f^{n})\}_{n\geq 1}$, we have
		\begin{align*}
			\sum_{i}\exp(\sup_{y\in B_{n_{i}}(x_{i},r)}-\tilde{\varphi}^{t}(y,f^{n_{i}})\leq \sum_{i}\exp(\sup_{z\in B_{n_{i}-m}(f^{m}(x_{i}),r)}-\tilde{\varphi}^{t}(z,f^{n_{i}-m})D,
		\end{align*}
		where $D=\sup_{x}-\tilde{\varphi}^{t}(x,f^{m})$.
		
		Fix such a number $r>0$, it follows from the definition of the topological pressure as the Carath$\acute{\text{e}}$odory singular dimension that
		$$m(\Lambda,\tilde{\Phi}_{f}(t),P(f|_\Lambda,\tilde{\Phi}_{f}(t))-\delta,r)=+\infty.$$
		Hence , for each $K>0$, there exists $L\in \mathbb{N}$ so that for any $N\geq L$, we have
		\begin{align*}
			&e^{-N(P(f|_{\Lambda},\tilde{\Phi}_{f}(t))-\delta)}\inf\{\sum_{i}\exp(\sup_{y\in B_{n_{i}}(x_{i},r)}-\tilde{\varphi}^{t}(y,f^{n_{i}})\}\\
			&\geq \inf\{\sum_{i}\exp(-(P(f|_{\Lambda},\tilde{\Phi}_{f}(t))-\delta)n_{i}+\sup_{y\in B_{n_{i}}(x_{i},r)}-\tilde{\varphi}^{t}(y,f^{n_{i}})\}\geq K.
		\end{align*}
		This yields that
		\begin{align*}
			\inf\{\sum_{i}\exp(\sup_{y\in B_{n_{i}}(x_{i},r)}-\tilde{\varphi}^{t}(y,f^{n_{i}})\}\geq e^{N(P(f|_{\Lambda},\tilde{\Phi}_{f}(t))-\delta)}K
		\end{align*}
		where the infimum is taken over all collections $\{B_{n_{i}}(x_{i},r)\}$ of Bowen's balls with $x_{i}\in \Lambda$, $n_{i}\geq N$, which cover $\Lambda$.
		
		Letting $N$ tends to infinity, we have that
		$$m(\Lambda\cap W^{u}(x),t,r)=+\infty$$
		for any $t<\alpha_{0}$. This implies that  $\dim_{C}\Lambda\cap W^{u}(x)> t$.
	\end{proof}
	
	
	\begin{Rem}
		As a result, we can use the notation $\dim_{C}\Lambda\cap W^{u}_{\text{loc}}(x)$ to express the Carath$\acute{\text{e}}$odory singular dimension of the set $\Lambda\cap W^{u}_{\text{loc}}(x)$.
	\end{Rem}
	\fi
	
	\iffalse
	\begin{Prop}\cite{caopesinzhao2019}\label{relation}
		Let $\Psi=\{\psi_{n}\}_{n\geq 1}$ be a super-additive sequence of continuous potentials on $X$. Then
		\begin{align*}
			P_{var}(f,\Psi)=\lim_{n\rightarrow \infty}P(f,\frac{\psi_{n}}{n})=\lim_{n\rightarrow\infty}\frac{1}{n}P(f^{n},\psi_{n}).
		\end{align*}
	\end{Prop}
	\begin{Lem}
		Assume that the map $f|_{\Lambda}$ possesses a $\{\lambda_{j}\}$-dominated splitting $T_{\Lambda}M=E_{1}\oplus E_{2}\oplus...\oplus E_{l}$ with $E_{1}\succeq E_{2}\succeq ...\succeq E_{k}$ and $\lambda_{1}>\lambda_{2}>...>\lambda_{\ell}>0>\lambda_{\ell+1}>...>\lambda_{k}$. Then for any $x\in \Lambda$, we have $dim_{H}\Lambda\cap W^{u}_{\text{loc}}(x)\geq s_{*}$, where $s_{*}$ is the unique root of Bowen's equation $\hat{P}_{sup}(s)=0$.
	\end{Lem}
	\begin{proof}
		Notice that for each $n\in \mathbb{N}$, $\Lambda$ is also a hyperbolic set for $f^{2^{n}}$. By Lemma~\ref{lower bound 1}, for each $n\in \mathbb{N}$, one has $dim_{H}\Lambda\geq s_{n}$, where $s_{n}$ is the unique root of the equation
		\begin{align*}
			P(f^{2^{n}}|_{\Lambda},-\hat{\psi}^{s}(\cdot,f^{2^{n}}))=0.
		\end{align*}
		By the variational principle of the topological pressure and the super-additivity of $\{-\hat{\psi}^{s}(\cdot,f^{n})\}_{n\geq 1}$, one can show that
		\begin{align*}
			\frac{1}{2^{k+1}}P(f^{2^{k+1}}|_{\Lambda},-\hat{\psi}^{s}(\cdot,f^{2^{k+1}}))\geq P(f^{2^{k}}|_{\Lambda},-\hat{\psi}^{s}(\cdot,f^{2^{k}})).
		\end{align*}
		Therefore, $s_{n}\leq s_{n+1}$. It follows that the limit $\lim_{n\rightarrow \infty}s_{n}$ exists and we denote it by $s_{*}$. Hence $dim_{H}\Lambda \cap W^{u}_{\text{loc}}(x)\geq s_{*}$. The results follows from Proposition ~\ref{relation}.
	\end{proof}
	\fi
	
	
	
	\section{Proof of the main results}\label{proof}
	In this section, we present the proof for our main results. %The proof is based on the lower bound for the repeller obtained by Cao, Pesin and Zhao in \cite{caopesinzhao2019}.
	\begin{Lem}\label{dim of measure}
		Suppose $f:M\rightarrow M$ is a $C^{1+\gamma}$ map on a $d-$dimensional Riemannian manifold and $\mu$ is an ergodic $f-$invariant measure which is absolutely continuous with the Lebesgue measure. Then $dim_{H} \mu=d$.
	\end{Lem}
	\begin{proof}
		For any measurable set $B$ with $\mu(B)=1$, by the absolute continuity of the measure $\mu$, we have $Leb(B)>0$. And hence $\dim_{H} B=d$. By the definition of $\dim_{H} \mu$, we get that $\dim_{H} \mu=d$.
	\end{proof}
	\begin{proof}[\textbf{Proof of Theorem~\ref{expanding}}]
		Since $f$ is expanding on $M$, there is an ergodic measure $\mu$ absolutely continuous with Lebesgue measure. Therefore, the Pesin's entropy formula ~\cite{BP} holds, that is
		\begin{align*}
			h_{\mu}(f)-\sum^{d}_{i=1} \lambda_{i}(\mu)=0.
		\end{align*}
		Hence, $$h_{\mu}(f)-\sum^{d-1}_{i=1} \lambda_{i}(\mu)=\lambda_{d}(\mu)>0.$$
		
		Fix $0<\varepsilon<\frac{\lambda_{d}(\mu)}{d}$ and $y\in M$, by Theorem~\ref{approximate}, there exists a repeller $\Lambda_{\varepsilon}$ satisfying $\Lambda_{\varepsilon}\cap \{y\}=\varnothing$ such that
		$$h_{\text{top}}(f|_{\Lambda_{\varepsilon}})\geq h_{\mu}(f)-\varepsilon$$
		and for any invariant ergodic measure $\nu$ supported on $\Lambda_{\varepsilon}$, we have
		$$|\lambda_{i}(\mu)-\lambda_{i}(\nu)|\leq \varepsilon\quad \text{for} \quad 1\leq i\leq d.$$
		Notice that $\Lambda_{\varepsilon}\cap \{y\}=\varnothing$ implies that $\Lambda_{\varepsilon}\subset E^{+}_{f}(y)$. Since $f$ is expanding on $\Lambda_{\varepsilon}$, there exists an invariant ergodic measure $\nu$ supported on $\Lambda_{\varepsilon}$ such that $$h_{\nu}(f|_{\Lambda_{\varepsilon}})=h_{\text{top}}(f|_{\Lambda_{\varepsilon}}).$$ Denote by $s^{\ast}$ the unique root of the equation $P_{\text{sup}}(s)=0$. By the definition of the super-additive pressure, we have
		\begin{equation}\label{entropy}
			h_{\nu}(f|_{\Lambda_{\varepsilon}})-\sum_{i=1}^{[s^{\ast}]}\lambda_{i}(\nu)-(s^{\ast}-[s^{\ast}])\lambda_{[s^{\ast}]+1}(\nu)\leq 0.
		\end{equation}
		By Theorem~\ref{approximate}, we have
		\begin{align*}
			\lambda_{d}(\mu)&=h_{\mu}(f)-\sum^{d-1}_{i=1} \lambda_{i}(\mu)\\
			&\leq h_{\nu}(f|_{\Lambda_{\varepsilon}})+\varepsilon-\sum_{i=1}^{d-1}\lambda_{i}(\nu)+(d-1)\varepsilon\\
			&<h_{\nu}(f|_{\Lambda_{\varepsilon}})-\sum_{i=1}^{d-1}\lambda_{i}(\nu)+\lambda_{d}(\mu).
		\end{align*}
		Thus, $$h_{\nu}(f|_{\Lambda_{\varepsilon}})-\sum_{i=1}^{d-1}\lambda_{i}(\nu)>0.$$ Combining this with (\ref{entropy}), we get that $s^{\ast}> d-1$. On the other hand, by Theorem~\ref{lower bound} and the monotonicity of the Hausdorff dimension, we have $$d\geq \dim_{H}E^{+}_{f}(y)\geq \dim_{H} \Lambda_{\varepsilon}\geq s^*.$$
		If $s^{\ast}=d$, there is nothing to prove. If $s^{\ast}<d$, then by Theorem~\ref{lower bound}, Theorem \ref{approximate} and (\ref{entropy}), we have
		\begin{align*}
			\dim_{H} \Lambda_{\varepsilon}&\geq s^{\ast}\\
			&\geq \frac{h_{\nu}(f|_{\Lambda_{\varepsilon}})-\sum_{i=1}^{[s^{\ast}]}\lambda_{i}(\nu)+[s^{\ast}]\lambda_{[s^{\ast}]+1}(\nu)}{\lambda_{[s^{\ast}]+1}(\nu)}\\
			&\geq \frac{h_{\mu}(f)-\varepsilon-\sum_{i=1}^{[s^{\ast}]}\lambda_{i}(\mu)-[s^{\ast}]\varepsilon+[s^{\ast}]\lambda_{[s^{\ast}]+1}(\mu)-[s^{\ast}]\varepsilon}{\lambda_{[s^{\ast}]+1}(\mu)+\varepsilon}\\
			&=\frac{([s^{\ast}]+1)\lambda_{[s^{\ast}]+1}(\mu)-(2[s^{\ast}]+1)\varepsilon}{\lambda_{[s^{\ast}]+1}(\mu)+\varepsilon}\\
			&=d-\frac{(3d-1)\varepsilon}{\lambda_{d}(\mu)+\varepsilon}.
		\end{align*}
		By the monotonicity of the Hausdorff dimension, we have $$d\geq \dim_{H}E^{+}_{f}(y)\geq \dim_{H} \Lambda_{\varepsilon}\geq d-\frac{(3d-1)\varepsilon}{\lambda_{d}(\mu)+\varepsilon}.$$ The arbitrariness of $\varepsilon>0$ yields that $\dim_{H}E^{+}_{f}(y)=d$. By Lemma \ref{dim of measure}, we finish the proof.
		
		
		
	\end{proof}
	%For the case that $f$ is not expanding on the whole manifold but admitting an repeller, we use the result in \cite{caopesinzhao2019} that the Carath$\acute{\text{e}}$odory singular dimension of the repeller is the unique root of the corresponding Bowen's equation to get our result.
	\begin{proof}[\textbf{Proof of Theorem~\ref{expanding2}}]
		Denote by $\alpha_{0}$ the unique root of $P(f|_{\Lambda},\alpha)=0$. By the variational principle of the sub-additive topological pressure and $f$ is expanding on $\Lambda$, there exists an ergodic measure $\mu^{\ast}$ such that
		\begin{align}\label{Pesin f1}
			h_{\mu^{\ast}}(f|_{\Lambda})-\sum^{d}_{i=d-[\alpha_{0}]+1}\lambda_{i}(\mu^{\ast})-(\alpha_{0}-[\alpha_{0}])\lambda_{d-[\alpha_{0}]}(\mu^{\ast})=0.
		\end{align}
		By Theorem~\ref{singular-valued dimension}, $\dim_{C} \Lambda=\alpha_{0}$.
		
		If $\alpha_{0}$ is an integer, then fix $0<\varepsilon<\frac{\lambda_{d-[\alpha_{0}]+1}(\mu^{\ast})}{[\alpha_{0}]}$, else fix $0<\varepsilon< \frac{(\alpha_{0}-[\alpha_{0}])\lambda_{d-[\alpha_{0}]}(\mu^{\ast})}{[\alpha_{0}]+1}$.
		For such an  $\varepsilon>0$ and any $y\in \Lambda$, by Theorem ~\ref{approximate}, there exists an $f-$invariant compact set $\Lambda_{\varepsilon}\subset \Lambda$ satisfying $\Lambda_{\varepsilon}\cap \{y\}=\varnothing$ such that $$h_{\text{top}}(f|_{\Lambda_{\varepsilon}})\geq h_{\mu^{\ast}}(f|_{\Lambda})-\varepsilon$$ and for any invariant measure $\nu$ supported on $\Lambda_{\varepsilon}$, we have $$|\lambda_{i}(\nu)-\lambda_{i}(\mu^{\ast})|\leq \varepsilon \quad \text{for} \quad 1\leq i\leq d.$$ Notice that $\Lambda_{\varepsilon}\subset E^{+}_{f|_{\Lambda}}(y)$. Since $f$ is expanding on $\Lambda_{\varepsilon}$, there exists an ergodic measure $\nu^{\ast}$ such that $$h_{\text{top}}(f|_{\Lambda_{\varepsilon}})=h_{\nu^{\ast}}(f|_{\Lambda_{\varepsilon}}).$$ Denote by $\alpha_{0}(\varepsilon)$ the unique root of $P(f|_{\Lambda_{\varepsilon}},\alpha)=0$.
		Then by the variational principle of the sub-additive topological pressure, we have
		\begin{equation}\label{entropy formula}
			h_{\nu^{\ast}}(f|_{\Lambda_{\varepsilon}})-\sum^{d}_{i=d-[\alpha_{0}(\varepsilon)]+1}\lambda_{i}(\nu^{\ast})-(\alpha_{0}(\varepsilon)-[\alpha_{0}(\varepsilon)])\lambda_{d-[\alpha_{0}(\varepsilon)]}(\nu^{\ast})\leq 0.
		\end{equation}
		If $\alpha_{0}$ is not an integer, then by Theorem \ref{approximate}, (\ref{Pesin f1}) and (\ref{entropy formula}), we have
		\begin{align*}
			(\alpha_{0}-[\alpha_{0}])\lambda_{d-[\alpha_{0}]}(\mu^{\ast})&=h_{\mu^{\ast}}(f|_{\Lambda})-\sum^{d}_{i=d-[\alpha_{0}]+1}\lambda_{i}(\mu^{\ast})\\
			&\leq h_{\nu^{\ast}}(f|_{\Lambda_{\varepsilon}})+\varepsilon-\sum^{d}_{i=d-[\alpha_{0}]+1}\lambda_{i}(\nu^{\ast})+[\alpha_{0}]\varepsilon\\
			&< h_{\nu^{\ast}}(f|_{\Lambda_{\varepsilon}})-\sum^{d}_{i=d-[\alpha_{0}]+1}\lambda_{i}(\nu^{\ast})+(\alpha_{0}-[\alpha_{0}])\lambda_{d-[\alpha_{0}]}(\mu^{\ast}).
		\end{align*}
		Combining this with (\ref{entropy formula}), we have \begin{align*}
			\sum^{d}_{i=d-[\alpha_{0}]+1}\lambda_{i}(\nu^{\ast})<\sum^{d}_{i=d-[\alpha_{0}(\varepsilon)]+1}\lambda_{i}(\nu^{\ast})+(\alpha_{0}(\varepsilon)-[\alpha_{0}(\varepsilon)])\lambda_{d-[\alpha_{0}(\varepsilon)]}(\nu^{\ast}).
		\end{align*}
		Therefore,
		$[\alpha_{0}]< \alpha_{0}(\varepsilon)$.
		On the other hand, by the monotonicity of the Carath$\acute{\text{e}}$odory singular dimension, we have $\dim_{C}\Lambda_{\varepsilon}\leq \dim_{C}\Lambda$, i.e. $\alpha_{0}(\varepsilon)\leq \alpha_{0}$. Hence, we have $[\alpha_{0}]=[\alpha_{0}(\varepsilon)]$.
		
		Therefore, by Theorem \ref{singular-valued dimension},  Theorem \ref{approximate}, (\ref{Pesin f1}) and (\ref{entropy formula}), we have
		\begin{align*}
			\text{dim}_{C}\Lambda_{\varepsilon}&=\alpha_{0}(\varepsilon)\geq\frac{h_{\nu^{\ast}}(f|_{\Lambda_{\varepsilon}})-\sum^{d}_{i=d-[\alpha_{0}(\varepsilon)]+1}\lambda_{i}(\nu^{\ast})+[\alpha_{0}(\varepsilon)]\lambda_{d-[\alpha_{0}(\varepsilon)]}(\nu^{\ast})}{\lambda_{d-[\alpha_{0}(\varepsilon)]}(\nu^{\ast})}\\
			&\geq\frac{h_{\mu^{\ast}}(f|_{\Lambda})-\varepsilon-\sum^{d}_{i=d-[\alpha_{0}]+1}\lambda_{i}(\mu^{\ast})-[\alpha_{0}]\varepsilon+[\alpha_{0}]\lambda_{d-[\alpha_{0}]}(\mu^{\ast})-[\alpha_{0}]\varepsilon}{\lambda_{d-[\alpha_{0}]}(\mu^{\ast})+\varepsilon}\\
			%&\geq \frac{h_{\mu^{\ast}}(f|_{\Lambda})-\sum^{d_{0}}_{i=d_{0}-[\alpha_{0}]+1}\lambda_{i}(\mu^{\ast})-(\alpha_{0}-[\alpha_{0}])\lambda_{d_{0}-[\alpha_{0}]}(\mu^{\ast})+\alpha_{0}\lambda_{d_{0}-[\alpha_{0}]}(\mu^{\ast})-(2[\alpha_{0}]+1)\varepsilon}{\lambda_{d_{0}-[\alpha_{0}]}(\mu^{\ast})+\varepsilon}\\
			&=\alpha_{0}-\frac{(\alpha_{0}+2[\alpha_{0}]+1)\varepsilon}{\lambda_{d-[\alpha_{0}]}(\mu^{\ast})+\varepsilon}.
		\end{align*}
		By the monotonicity of the Carath$\acute{\text{e}}$odory singular dimension, we have $$\dim_{C}\Lambda\geq \dim_{C} E^{+}_{f|_{\Lambda}}(y)\geq \dim_{C}\Lambda_{\varepsilon}\geq \alpha_{0}-\frac{(\alpha_{0}+2[\alpha_{0}]+1)\varepsilon}{\lambda_{d-[\alpha_{0}]}(\mu^{\ast})+\varepsilon}.$$ Let $\varepsilon$ tend to $0$. This yields the result.
		
		If $\alpha_{0}$ is  an integer, then by Theorem \ref{approximate}, (\ref{Pesin f1}) and (\ref{entropy formula}), we have
		\begin{align*}
			\lambda_{d-[\alpha_{0}]+1}(\mu^{\ast})&=h_{\mu^{\ast}}(f|_{\Lambda})-\sum^{d}_{i=d-[\alpha_{0}]+2}\lambda_{i}(\mu^{\ast})\\
			&\leq h_{\nu^{\ast}}(f|_{\Lambda_{\varepsilon}})+\varepsilon-\sum^{d}_{i=d-[\alpha_{0}]+2}\lambda_{i}(\nu^{\ast})+([\alpha_{0}]-1)\varepsilon\\
			&< h_{\nu^{\ast}}(f|_{\Lambda_{\varepsilon}})-\sum^{d}_{i=d-[\alpha_{0}]+2}\lambda_{i}(\nu^{\ast})+\lambda_{d-[\alpha_{0}]+1}(\mu^{\ast}).
		\end{align*}
		Combining this with ~(\ref{entropy formula}), we have 
		\begin{align*}
			\sum^{d}_{i=d-[\alpha_{0}]+2}\lambda_{i}(\nu^{\ast})<\sum^{d}_{i=d-[\alpha_{0}(\varepsilon)]+1}\lambda_{i}(\nu^{\ast})+(\alpha_{0}(\varepsilon)-[\alpha_{0}(\varepsilon)])\lambda_{d-[\alpha_{0}(\varepsilon)]}(\nu^{\ast}).
		\end{align*}
		
		Therefore, $[\alpha_{0}]-1< \alpha_{0}(\varepsilon)$. Since $\alpha_0(\varepsilon)\leq \alpha_0$ and $\alpha_{0}$ is  an integer, we have $\alpha_0-1<\alpha_0(\varepsilon)\leq \alpha_0$. Hence $\alpha=\alpha_{0}(\varepsilon)$ or $[\alpha_{0}]= [\alpha_{0}(\varepsilon)]+1$.
		
		If $\alpha_{0}=\alpha_{0}(\varepsilon)$, then the result follows from the monotonicity of the Carath$\acute{\text{e}}$odory singular dimension.
		
		If $[\alpha_{0}]=[\alpha_{0}(\varepsilon)]+1$, then by Theorem \ref{singular-valued dimension},  Theorem \ref{approximate}, (\ref{Pesin f1}) and (\ref{entropy formula}), we have
		\begin{align*}
			&\dim_{C}\Lambda_{\varepsilon}=\alpha_{0}(\varepsilon)\geq\frac{h_{\nu^{\ast}}(f|_{\Lambda_{\varepsilon}})-\sum^{d}_{i=d-[\alpha_{0}(\varepsilon)]+1}\lambda_{i}(\nu^{\ast})+[\alpha_{0}(\varepsilon)]\lambda_{d-[\alpha_{0}(\varepsilon)]}(\nu^{\ast})}{\lambda_{d-[\alpha_{0}(\varepsilon)]}(\nu^{\ast})}\\
			&\geq\frac{h_{\mu^{\ast}}(f|_{\Lambda})-\varepsilon-\sum^{d}_{i=d-[\alpha_{0}(\varepsilon)]+1}\lambda_{i}(\mu^{\ast})-[\alpha_{0}(\varepsilon)]\varepsilon+[\alpha_{0}(\varepsilon)]\lambda_{d-[\alpha_{0}(\varepsilon)]}(\mu^{\ast})-[\alpha_{0}(\varepsilon)]\varepsilon}{\lambda_{d-[\alpha_{0}(\varepsilon)]}(\mu^{\ast})+\varepsilon}\\
			&=\frac{h_{\mu^{\ast}}(f|_{\Lambda})-\sum^{d}_{i=d-[\alpha_{0}]+1}\lambda_{i}(\mu^{\ast})-(2[\alpha_{0}]-1)\varepsilon+[\alpha_{0}]\lambda_{d-[\alpha_{0}]+1}(\mu^{\ast})}{\lambda_{d-[\alpha_{0}]+1}(\mu^{\ast})+\varepsilon}\\ &=\alpha_{0}-\frac{(3[\alpha_{0}]-1)\varepsilon}{\lambda_{d-[\alpha_{0}]+1}(\mu^{\ast})+\varepsilon}.
		\end{align*}
		
		The monotonicity of the Carath$\acute{\text{e}}$odory singular dimension and the arbitrariness of  $\varepsilon>0$ imply the result.
	\end{proof}
	\iffalse
	\begin{Lem}
		Suppose $\Lambda$ is a hyperbolic attractor for a $C^{1+\alpha}$ diffeomorphism $f$ and $\mu$ is the SRB measure for $f$. Then for any $x\in \Lambda$, we have $dim_{H}\mu_{x}^{u}=dim W^{u}(x)$.
	\end{Lem}
	\begin{proof}
		For any measurable set $B$ with $\mu_{x}^{u}(B)=1$, we have $Leb^{u}(B)>0$. And hence $\dim_{H} B=\dim W^{u}(x)$. By the definition of $\dim_{H}\mu_{x}^{u}$, we get the result.
	\end{proof}
	
	
	\begin{proof}[\textbf{Proof of Theorem~\ref{hyperbolic1}}]
		Since $M$ is a transitive hyperbolic set for a $C^{1+\alpha}$ diffeomorphism $f$, there is a unique SRB measure $\mu$ for $f$. Therefore, the Pesin' entropy formula ~\cite{LS} holds, that is
		\begin{align*}
			h_{\mu}(f)-\sum^{u}_{i=1} \lambda_{i}(\mu)=0,
		\end{align*}
		here $u=\dim W^{u}(x)$.
		Hence, $$h_{\mu}(f)-\sum^{u-1}_{i=1} \lambda_{i}(\mu)=\lambda_{u}(\mu)>0.$$
		Fix $0<\varepsilon<\frac{\lambda_{u}(\mu)}{u}$ and $x,y\in M$, by Theorem~\ref{Katok}, there exists a hyperbolic set $\Lambda_{\varepsilon}$ satisfying $\Lambda_{\varepsilon}\cap \{y\}=\varnothing$ such that
		$$h_{\text{top}}(f,\Lambda_{\varepsilon})\geq h_{\mu}(f)-\varepsilon,$$
		and for any invariant measure $\nu$ supported on $\Lambda_{\varepsilon}$, we have $$|\lambda_{i}(\mu)-\lambda_{i}(\nu)|\leq \varepsilon\quad \text{for} \quad 1\leq i\leq u.$$ Notice that $\Lambda_{\varepsilon}\cap \{y\}=\varnothing$ implies that $\Lambda_{\varepsilon}\subset E^{+}_{f}(y)$. Since $f$ is hyperbolic on $\Lambda_{\varepsilon}$, there exists an invariant measure $\nu$ supported on $\Lambda_{\varepsilon}$ such that $$h_{\nu}(f|_{\Lambda_{\varepsilon}})=h_{\text{top}}(f,\Lambda_{\varepsilon}).$$ Denote by $s_{\ast}$ the unique root of the equation $P_{\text{var}}(f|_{\Lambda},\tilde{\Psi}_{f}(s))=0$. By the definition of the super-additive pressure, we have
		\begin{equation}\label{entropy2}
			h_{\nu}(f|_{\Lambda_{\varepsilon}})-\sum_{i=1}^{[s_{\ast}]}\lambda_{i}(\nu)-(s_{\ast}-[s_{\ast}])\lambda_{[s_{\ast}]+1}(\nu)\leq 0.
		\end{equation}
		By Theorem~\ref{Katok}, we have
		\begin{align*}
			\lambda_{u}(\mu)&=h_{\mu}(f)-\sum^{u-1}_{i=1} \lambda_{i}(\mu)\\
			&\leq h_{\nu}(f|_{\Lambda_{\varepsilon}})+\varepsilon-\sum_{i=1}^{u-1}\lambda_{i}(\nu)+(u-1)\varepsilon\\
			&<h_{\nu}(f|_{\Lambda_{\varepsilon}})-\sum_{i=1}^{u-1}\lambda_{i}(\nu)+\lambda_{u}(\mu).
		\end{align*}
		Thus, $$h_{\nu}(f|_{\Lambda_{\varepsilon}})-\sum_{i=1}^{u-1}\lambda_{i}(\nu)>0.$$ Combining this with (\ref{entropy2}), we get that $s_{\ast}> u-1$. On the other hand, by Theorem~\ref{lower bound for hyperbolic} and the monotonicity of the Hausdorff dimension, we have $s_{\ast}\leq \dim_{H}\Lambda_{\varepsilon}\cap W^{u}(x)\leq u$.
		
		If $s_{\ast}<u$, then by Theorem~\ref{lower bound for hyperbolic}, we have
		\begin{align*}
			\dim_{H} \Lambda_{\varepsilon}\cap W^{u}(x)&\geq s_{\ast}\\
			&\geq \frac{h_{\nu}(f|_{\Lambda_{\varepsilon}})-\sum_{i=1}^{[s_{\ast}]}\lambda_{i}(\nu)+[s_{\ast}]\lambda_{[s_{\ast}]+1}(\nu)}{\lambda_{[s_{\ast}]+1}(\nu)}\\
			&\geq \frac{h_{\mu}(f)-\varepsilon-\sum_{i=1}^{[s_{\ast}]}\lambda_{i}(\mu)-[s_{\ast}]\varepsilon+[s_{\ast}]\lambda_{[s_{\ast}]+1}(\mu)-[s_{\ast}]\varepsilon}{\lambda_{[s_{\ast}]+1}(\mu)+\varepsilon}\\
			&=\frac{([s_{\ast}]+1)\lambda_{[s_{\ast}]+1}(\mu)-(2[s_{\ast}]+1)\varepsilon}{\lambda_{[s_{\ast}]+1}(\mu)+\varepsilon}\\
			&=u-\frac{(3u-1)\varepsilon}{\lambda_{u}(\mu)+\varepsilon}.
		\end{align*}
		By the monotonicity of the Hausdorff dimension, we have $$u\geq \dim_{H}E^{+}_{f}(y)\cap W^{u}(x)\geq \dim_{H} \Lambda_{\varepsilon}\cap W^{u}(x)\geq u-\frac{(3u-1)\varepsilon}{\lambda_{u}(\mu)+\varepsilon}.$$ The arbitrariness of  the $\varepsilon>0$ implies that $\dim_{H}E^{+}_{f}(y)\cap W^{u}(x)=u$.
		
		For the case that $s_{\ast}=u$, the proof is similar, so we omit it.
	\end{proof}
	
	
	
	
	\begin{proof}[\textbf{Proof of Theorem~\ref{hyperbolic2}}]
		Denote by $\alpha_{0}$ the unique root of $P(f|_{\Lambda},\tilde{\Phi}_{f}(\alpha))=0$. By the variational principle of the sub-additive topological pressure and $f$ is hyperbolic, there exists an ergodic measure $\mu^{\ast}$ such that
		\begin{align*}
			h_{\mu^{\ast}}(f|_{\Lambda})-\sum^{u}_{i=u-[\alpha_{0}]+1}\lambda_{i}(\mu^{\ast})-(\alpha_{0}-[\alpha_{0}])\lambda_{u-[\alpha_{0}]}(\mu^{\ast})=0.
		\end{align*}
		By Theorem~\ref{singular-valued dimension for hyperbolic}, for any $x\in\Lambda$, we have $\dim_{C} \Lambda\cap W^{u}_{\text{loc}}(x)=\alpha_{0}$.
		
		If $\alpha_{0}$ is an integer, then fix $0<\varepsilon<\frac{\lambda_{u-[\alpha_{0}]+1}(\mu^{\ast})}{[\alpha_{0}]}$, else fix $0<\varepsilon< \frac{(\alpha_{0}-[\alpha_{0}])\lambda_{u-[\alpha_{0}]}(\mu^{\ast})}{[\alpha_{0}]+1}$.
		For such an  $\varepsilon>0$ and any $y\in \Lambda$, by Theorem ~\ref{Katok}, there exists an $f-$invariant compact set $\Lambda_{\varepsilon}\subset \Lambda$ satisfying $\Lambda_{\varepsilon}\cap \{y\}=\varnothing$ such that $$h_{\text{top}}(f,\Lambda_{\varepsilon})\geq h_{\mu^{\ast}}(f|_{\Lambda})-\varepsilon$$ and for any invariant measure $\nu$ supported on $\Lambda_{\varepsilon}$, we have $$|\lambda_{i}(\nu)-\lambda_{i}(\mu^{\ast})|\leq \varepsilon \quad \text{for}\quad 1\leq i\leq u.$$ Notice that $\Lambda_{\varepsilon}\subset E^{+}_{f|_{\Lambda}}(y)$.
		
		Since $f$ is hyperbolic on $\Lambda_{\varepsilon}$, there exists an ergodic measure $\nu^{\ast}$ such that $$h_{\text{top}}(f,\Lambda_{\varepsilon})=h_{\nu^{\ast}}(f|_{\Lambda_{\varepsilon}}).$$ Denote by $\alpha_{0}(\varepsilon)$ the unique root of $P(f|_{\Lambda_{\varepsilon}},\tilde{\Phi}_{f}(\alpha))=0$.
		Then by the variational principle of the sub-additive topological pressure, we have
		\begin{equation}\label{entropy formula2}
			h_{\nu^{\ast}}(f|_{\Lambda_{\varepsilon}})-\sum^{u}_{i=u-[\alpha_{0}(\varepsilon)]+1}\lambda_{i}(\nu^{\ast})-(\alpha_{0}(\varepsilon)-[\alpha_{0}(\varepsilon)])\lambda_{u-[\alpha_{0}(\varepsilon)]}(\nu^{\ast})\leq 0.
		\end{equation}
		If $\alpha_{0}$ is not an integer, then
		\begin{align*}
			(\alpha_{0}-[\alpha_{0}])\lambda_{u-[\alpha_{0}]}(\mu^{\ast})&=h_{\mu^{\ast}}(f|_{\Lambda})-\sum^{u}_{i=u-[\alpha_{0}]+1}\lambda_{i}(\mu^{\ast})\\
			&\leq h_{\nu^{\ast}}(f|_{\Lambda_{\varepsilon}})+\varepsilon-\sum^{u}_{i=u-[\alpha_{0}]+1}\lambda_{i}(\nu^{\ast})+[\alpha_{0}]\varepsilon\\
			&< h_{\nu^{\ast}}(f|_{\Lambda_{\varepsilon}})-\sum^{u}_{i=u-[\alpha_{0}]+1}\lambda_{i}(\nu^{\ast})+(\alpha_{0}-[\alpha_{0}])\lambda_{u-[\alpha_{0}]}(\mu^{\ast}).
		\end{align*}
		Combining this with (\ref{entropy formula2}), we have $[\alpha_{0}]\leq [\alpha_{0}(\varepsilon)]$ and $\alpha_{0}(\varepsilon)$ is also not an integer.
		On the other hand, by the monotonicity of the Carath$\acute{\text{e}}$odory singular dimension, we have $\dim_{C}\Lambda_{\varepsilon}\cap W^{u}(x)\leq \dim_{C}\Lambda\cap W^{u}(x)$, i.e. $\alpha_{0}(\varepsilon)\leq \alpha_{0}$. Hence, we have $[\alpha_{0}]=[\alpha_{0}(\varepsilon)]$.
		
		Therefore,
		\begin{align*}
			&\dim_{C}\Lambda_{\varepsilon}=\alpha_{0}(\varepsilon)\geq\frac{h_{\nu^{\ast}}(f|_{\Lambda_{\varepsilon}})-\sum^{u}_{i=u-[\alpha_{0}(\varepsilon)]+1}\lambda_{i}(\nu^{\ast})+[\alpha_{0}(\varepsilon)]\lambda_{u-[\alpha_{0}(\varepsilon)]}(\nu^{\ast})}{\lambda_{u-[\alpha_{0}(\varepsilon)]}(\nu^{\ast})}\\
			&\geq\frac{h_{\mu^{\ast}}(f|_{\Lambda})-\varepsilon-\sum^{u}_{i=u-[\alpha_{0}]+1}\lambda_{i}(\mu^{\ast})-[\alpha_{0}]\varepsilon+[\alpha_{0}]\lambda_{u-[\alpha_{0}]}(\mu^{\ast})-[\alpha_{0}]\varepsilon}{\lambda_{u-[\alpha_{0}]}(\mu^{\ast})+\varepsilon}\\
			&\geq \frac{h_{\mu^{\ast}}(f|_{\Lambda})-\sum^{u}_{i=u-[\alpha_{0}]+1}\lambda_{i}(\mu^{\ast})-(\alpha_{0}-[\alpha_{0}])\lambda_{u-[\alpha_{0}]}(\mu^{\ast})+\alpha_{0}\lambda_{u-[\alpha_{0}]}(\mu^{\ast})-(2[\alpha_{0}]+1)\varepsilon}{\lambda_{u-[\alpha_{0}]}(\mu^{\ast})+\varepsilon}\\
			&=\alpha_{0}-\frac{(\alpha_{0}+2[\alpha_{0}]+1)\varepsilon}{\lambda_{d-[\alpha_{0}]}(\mu^{\ast})+\varepsilon}.
		\end{align*}
		By the monotonicity of the Carath$\acute{\text{e}}$odory singular dimension, we have $$\dim_{C}\Lambda\cap W^{u}(x)\geq \dim_{C} E^{+}_{f|_{\Lambda}}(y)\cap W^{u}(x)\geq \dim_{C}\Lambda_{\varepsilon}\cap W^{u}(x)\geq \alpha_{0}-\frac{(\alpha_{0}+2[\alpha_{0}]+1)\varepsilon}{\lambda_{u-[\alpha_{0}]}(\mu^{\ast})+\varepsilon}.$$ Let $\varepsilon$ tend to $0$. This yields the result.
		
		If $\alpha_{0}$ is  an integer, then
		\begin{align*}
			\lambda_{u-[\alpha_{0}]+1}(\mu^{\ast})&=h_{\mu^{\ast}}(f|_{\Lambda})-\sum^{u}_{i=u-[\alpha_{0}]+2}\lambda_{i}(\mu^{\ast})\\
			&\leq h_{\nu^{\ast}}(f|_{\Lambda_{\varepsilon}})+\varepsilon-\sum^{u}_{i=u-[\alpha_{0}]+2}\lambda_{i}(\nu^{\ast})+([\alpha_{0}]-1)\varepsilon\\
			&< h_{\nu^{\ast}}(f|_{\Lambda_{\varepsilon}})-\sum^{u}_{i=u-[\alpha_{0}]+2}\lambda_{i}(\nu^{\ast})+\lambda_{u-[\alpha_{0}]+1}(\mu^{\ast}).
		\end{align*}
		Combining this with ~(\ref{entropy formula2}), we get that (i) $[\alpha_{0}]\leq [\alpha_{0}(\varepsilon)]$ and $\alpha_{0}(\varepsilon)$ is  an integer or (ii) $[\alpha_{0}]\leq [\alpha_{0}(\varepsilon)]+1$ and $\alpha_{0}(\varepsilon)$ is not an integer.
		
		If $[\alpha_{0}]\leq [\alpha_{0}(\varepsilon)]$, then this results that $\alpha_{0}=\alpha_{0}(\varepsilon)$.
		
		If $[\alpha_{0}]=[\alpha_{0}(\varepsilon)]+1$,  then
		\begin{align*}
			&\dim_{C}\Lambda_{\varepsilon}\cap W^{u}(x)=\alpha_{0}(\varepsilon)\geq\frac{h_{\nu^{\ast}}(f|_{\Lambda_{\varepsilon}})-\sum^{u}_{i=u-[\alpha_{0}(\varepsilon)]+1}\lambda_{i}(\nu^{\ast})+[\alpha_{0}(\varepsilon)]\lambda_{u-[\alpha_{0}(\varepsilon)]}(\nu^{\ast})}{\lambda_{u-[\alpha_{0}(\varepsilon)]}(\nu^{\ast})}\\
			&\geq\frac{h_{\mu^{\ast}}(f|_{\Lambda})-\varepsilon-\sum^{u}_{i=u-[\alpha_{0}(\varepsilon)]+1}\lambda_{i}(\mu^{\ast})-[\alpha_{0}(\varepsilon)]\varepsilon+[\alpha_{0}(\varepsilon)]\lambda_{u-[\alpha_{0}(\varepsilon)]}(\mu^{\ast})-[\alpha_{0}(\varepsilon)]\varepsilon}{\lambda_{u-[\alpha_{0}(\varepsilon)]}(\mu^{\ast})+\varepsilon}\\
			&=\frac{h_{\mu^{\ast}}(f|_{\Lambda})-\sum^{u}_{i=u-[\alpha_{0}]+1}\lambda_{i}(\mu^{\ast})-(2[\alpha_{0}]-1)\varepsilon+[\alpha_{0}]\lambda_{u-[\alpha_{0}]+1}(\mu^{\ast})}{\lambda_{u-[\alpha_{0}]+1}(\mu^{\ast})+\varepsilon}\\ &=\alpha_{0}-\frac{(3[\alpha_{0}]-1)\varepsilon}{\lambda_{u-[\alpha_{0}]+1}(\mu^{\ast})+\varepsilon}.
		\end{align*}
		
		
		
		By the monotonicity of the Carath$\acute{\text{e}}$odory singular dimension and then let $\varepsilon$ tend to $0$, this yields the result.
	\end{proof}
	\fi
	
	
	\begin{thebibliography}{Qu}
		\bibitem[ACW]{ACW}A. Avila, S. Crovisier and A. Wilkinson. $C^1$  density of stable ergodicity. Adv. Math., 379(5): 107496, 68 pp, 2021.
		\bibitem[BP]{BP} L. Barreira and Y. Pesin. Nonuniform hyperbolicity, Encyclopedia of Mathematics
		and its Applications. vol. 115, Cambridge University Press, Cambridge, 2007.
		\bibitem[CFH]{cfh}Y. Cao, D. Feng and W. Huang. The thermodynamic formalism for sub-additive potentials. Discrete Contin. Dyn. Syst., 20: 639-657, 2008.
		\bibitem[CG1]{CG1}S. Campos and K. Gelfert. Exceptional sets for nonuniformly expanding maps. Nonlinearity, 29:  1238-1256, 2016.
		\bibitem[CG2]{CG2}S. Campos and K. Gelfert. Exceptional sets for nonuniformly hyperbolic diffeomorphisms. J. Dynam. Differential Equations, 31(2): 979-1004, 2019.
		\bibitem[CliPZ]{CliPZ} V. Climenhaga, Y. Pesin and A. Zelerowicz. Equilibrium states in dynamical systems via geometric measure theory. Bull.
		Amer. Math. Soc. (N.S.), 56(4): 569-610, 2019.
		\bibitem[CPZ]{caopesinzhao2019}  Y. Cao, Y. Pesin and Y. Zhao. Dimension estimates for non-conformal repellers and continuity of sub-additive topological pressure. Geom. Funct. Anal., 29(5): 1325-1368, 2019. 
		\bibitem[Do]{Do}D. Dolgopyat. Bounded orbits of Anosov flows. Duke Math. J., 87:  87-114, 1997.
		\bibitem[G1]{gelfert2010}K. Gelfert. Repellers for non-uniformly expanding maps with singular or critical points. Bull.
		Braz. Math. Soc., 41: 237-257, 2010.
		\bibitem[Kat80]{katok1980} A. Katok. Lyapunov exponents, entropy and periodic orbits for diffeomorphisms. Publications Math$\acute{\text{e}}$matiques De Linstitut Des Hautes $\acute{\text{e}}$tudes Scientifiques, 51: 137-173, 1980.
		\bibitem[KH]{KH}A. Katok and B. Hasselblatt. Introduction to the Modern Theory of Dynamical Systems. Encyclopedia of Mathematics and Its Applications. vol. 54, Cambridge University Press, Cambridge, 1995.
		\bibitem[KW]{KW}D. Y. Kleinbock and B. Weiss. Modified Schmidt games and Diophantine approximation with weights. Adv. Math., 223(4): 1276-1298, 2010.
		\bibitem[LS]{LS}F. Ledrappier and J. Strelcyn. A proof of the estimation from below in Pesin's entropy formula. Ergod. Th. \& Dynam. Sys., 2: 203-219, 1983.
		%\bibitem[O]{Oseledets}V. Oseledets. A multiplicative ergodic theorem. Trans. Moscow. Math. Soc., 19:  197-231, 1968.
		
		\bibitem[Pes]{Pesin}Y. Pesin. Dimension Theory in Dynamical Systems. Contemporary Views and Applications. Chicago Lectures in Mathematics, University of Chicago Press, Chicago, IL, 1997.
		\bibitem[QW]{QW}C. Qu and J. Wang. Exceptional sets for average conformal dynamical systems. Chaos, Solitons \& Fractals, 161: 112351, 2022.
		\bibitem[Sch]{Schmidt}W. M. Schmidt. On badly approximable numbers and certain games. Trans. Amer. Math. Soc., 123(1): 178-199, 1966.
		\bibitem[Tse1]{Tseng1}J. Tseng. Schmidt games and Markov partitions. Nonlinearity, 22:  525-543, 2009.
		\bibitem[Tse2]{Tseng2}J. Tseng. Nondense orbits for Anosov diffeomorphisms of the 2-torus. Real Analysis Exchange, 41(2):  307-314, 2016.
		\bibitem[U]{Urbanski}M. Urba$\acute{\text{n}}$ski. The Hausdorff dimension of the set of points with nondense orbit under a hyperbolic dynamical system. Nonlinearity, 4: 385-397, 1991.
		
		\bibitem[Wu1]{Wu1}W. Wu. Schmidt games and non-dense forward orbits of certain partially hyperbolic systems. Ergod. Th. \& Dynam. Sys., 36(5): 1656-1678, 2016.
		\bibitem[Wu2]{Wu2}W. Wu. Modified Schmidt games and non-dense forward orbits of partially hyperbolic systems. Discrete Contin. Dyn. Syst. Ser. A, 36(6): 3463-3481, 2016.
		\bibitem[Wu3]{Wu3}W. Wu. On non-dense orbits of certain non-algebraic dynamical systems. J. Dynam. Differential Equations, 30(2): 501-519, 2018.
		\bibitem[Wal]{WALTERS} P. Walters. An Introduction to Ergodic Theory. Springer-Verlag, New York, 1982.
		\bibitem[You]{Young}L. S. Young. Dimension, entropy and Lyapunov exponents. Ergod. Th. \& Dynam. Sys., 2: 109-124, 1982.
	\end{thebibliography}
	
	
	
\end{document}
